% Classification des éléments — PCT 3ème — Cours signé OnBuch+
% Programme officiel MINESEC. Compiler avec : tectonic N02-classification-des-elements.tex
\newcommand{\DOCMATIERE}{PCT}
\newcommand{\DOCNIVEAU}{3ème}
\newcommand{\DOCMODULE}{Module 1 : Chimie}
\newcommand{\DOCLECON}{N02}
\newcommand{\DOCTITRE}{Classification des éléments}
% OnBuch+ — préambule « cours signés » ENRICHI (compilé avec tectonic / XeLaTeX)
% Système visuel « Pop » d'OnBuch+ : fond crème (jamais blanc), bordures encre
% épaisses, ombres « dures » décalées, titres Archivo Black en capitales.
% Version MATHÉMATIQUES : graphes (pgfplots/tikz), boîtes pédagogiques
% (définition, propriété, méthode, attention, le savais-tu, exercice résolu,
% exercice, TP, bilan), page de garde et signature OnBuch+ ancrée en bas.
%
% Macros attendues dans main.tex AVANT \input{preamble} :
%   \DOCMATIERE  \DOCNIVEAU  \DOCMODULE  \DOCLECON (numéro)  \DOCTITRE
\documentclass[11pt]{article}

\usepackage{fontspec}
\usepackage[french]{babel}
\usepackage[a4paper,margin=2cm,top=3.4cm,bottom=2.8cm,headheight=1.4cm,footskip=1.3cm]{geometry}
\usepackage{amsmath,amssymb,amsfonts}
\usepackage{mathtools} % \xmapsto, \coloneqq... souvent utilisés par les modèles
\usepackage{cancel} % simplifications barrées (bilans de forces)
% \abs{z} (module, valeur absolue) et \norm{u} (norme), utilisés par les modèles
\providecommand{\abs}{}\let\abs\relax\DeclarePairedDelimiter{\abs}{\lvert}{\rvert}
\providecommand{\norm}{}\let\norm\relax\DeclarePairedDelimiter{\norm}{\lVert}{\rVert}
\usepackage[table]{xcolor} % [table] : fournit aussi \rowcolors (alternance de lignes)
\usepackage{pagecolor}
\usepackage{tikz}
\usetikzlibrary{arrows.meta,positioning,calc,shapes.geometric,shapes.misc,shapes.arrows,decorations.pathmorphing,decorations.pathreplacing,decorations.markings,patterns,shadows,mindmap,trees,fit,backgrounds,matrix,intersections,angles,quotes}
\usepackage{pgfplots}
\usepackage[version=4]{mhchem} % équations nucléaires \ce{^{235}_{92}U}
\usepackage[european,siunitx]{circuitikz} % schémas électriques
\pgfplotsset{compat=1.18}
\usepgfplotslibrary{fillbetween,groupplots} % aires entre courbes (calcul intégral)
\usepackage{enumitem}
\usepackage{fancyhdr}
% [most] charge toutes les bibliothèques tcolorbox (listings, minted, raster,
% documentation...) et gonfle inutilement la mémoire TeX sur un document long
% (~300 boîtes) : on ne charge que ce qui est réellement utilisé.
\usepackage[skins,breakable]{tcolorbox}
\usepackage{graphicx}
\usepackage{colortbl}
\usepackage{booktabs}
\usepackage{tabularx}
\usepackage{multirow}
\usepackage{diagbox} % case d'en-tête diagonale des tableaux à double entrée
% Colonnes extensibles centrée / gauche / droite (C, L, R), souvent utilisées par les modèles
\newcolumntype{C}{>{\centering\arraybackslash}X}
\newcolumntype{L}{>{\raggedright\arraybackslash}X}
\newcolumntype{R}{>{\raggedleft\arraybackslash}X}
\usepackage{array}
\usepackage{multicol}
\usepackage{siunitx}
% parse-numbers=false : accepte aussi \SI{2,5 \times 10^{-3}}{...}
\sisetup{locale=FR,output-decimal-marker={,},group-minimum-digits=4,parse-numbers=false}
\DeclareSIUnit\ohm{\ensuremath{\symup{\Omega}}} % U+2126 absent de la police
\DeclareSIUnit\division{div} % réglages d'oscilloscope (ms/div, V/div)
\DeclareSIUnit\tr{tr} % tours (vitesse de rotation, ex. \SI{1500}{\tr\per\minute})
\DeclareSIUnit\molar{M} % concentration molaire (ex. \SI{3}{\molar}, \SI{1}{\micro\molar})
\DeclareSIUnit\an{an} % année (le modèle confond parfois avec \year, un compteur TeX) — ex. \SI{5}{\kilo\meter\per\an}
\DeclareSIUnit\FCFA{FCFA} % franc CFA (ex. \SI{500}{\FCFA\per\kilo\gram})
\DeclareSIUnit\degreeBrix{\textdegree Brix} % teneur en sucre d'un sirop/jus (réfractométrie)
\DeclareSIUnit\month{mois} % le modèle utilise parfois \month, un compteur TeX, comme unité de durée
\DeclareSIUnit\microliter{\micro\liter} % le modèle écrit parfois \microliter en un seul token
\DeclareSIUnit\million{millions} % dénombrement (ex. \SI{15}{\million\per\mL} spermatozoïdes)
\usepackage{qrcode}
% \degree : commande fréquemment utilisée par les modèles pour un angle ou une
% température en dehors de \SI{}{} (non fournie par les paquets chargés).
\providecommand{\degree}{^{\circ}}
% \qed (fin de démonstration), utilisé par les modèles sans amsthm
\providecommand{\qed}{\hfill$\square$}
% \barre : double barre des valeurs interdites dans un tableau de variations (tkz-tab)
\providecommand{\barre}{\ensuremath{\|}}
% environnement proof (amsthm non chargé) utilisé par les modèles
\ifdefined\proof\else\newenvironment{proof}[1][Démonstration]{\par\noindent\textit{#1.}\ }{\qed\par}\fi
\usepackage{lastpage}
\usepackage{hyperref}
\hypersetup{hidelinks}

% ============================================================
% Palette « Pop » OnBuch+ — mêmes hex que la classe P (plus_ui.dart)
% ============================================================
\definecolor{poporange}{HTML}{F59321}
\definecolor{popdark}{HTML}{E07A0C}
\definecolor{popcream}{HTML}{FFF6E8}
\definecolor{popink}{HTML}{1C1714}
\definecolor{poppurple}{HTML}{7A5AE0}
\definecolor{popgreen}{HTML}{2E9E6A}
\definecolor{poppink}{HTML}{E0457B}
\definecolor{popblue}{HTML}{2F7DE1}
\definecolor{popgold}{HTML}{C98A12}
\definecolor{popmuted}{HTML}{8A7F76}
\definecolor{popline}{HTML}{EADFD2}
\definecolor{popgray}{HTML}{8A7F76}
\colorlet{popgrey}{popgray}
% Alias tolérants (le modèle nomme parfois les couleurs différemment)
\colorlet{popwhite}{white}
\colorlet{popblack}{popink}
\colorlet{popred}{poppink}
\colorlet{popyellow}{popgold}
% Teintes claires (fonds de tableaux/encadrés) — le modèle nomme parfois
% les couleurs avec un suffixe L (ex. popblueL) sans les avoir définies.
\colorlet{poporangeL}{poporange!20}
\colorlet{popdarkL}{popdark!20}
\colorlet{poppurpleL}{poppurple!20}
\colorlet{popgreenL}{popgreen!20}
\colorlet{poppinkL}{poppink!20}
\colorlet{popblueL}{popblue!20}
\colorlet{popgoldL}{popgold!20}
\colorlet{popredL}{popred!20}
\colorlet{popgrayL}{popgray!20}
% Teintes claires pour fonds de tableaux / zones
\colorlet{popblueL}{popblue!10!white}
\colorlet{poporangeL}{poporange!14!white}
\colorlet{popgreenL}{popgreen!12!white}
\colorlet{poppurpleL}{poppurple!12!white}
\colorlet{poppinkL}{poppink!12!white}
\colorlet{popgoldL}{popgold!14!white}

\pagecolor{popcream}

% ============================================================
% Typographie
% ============================================================
\defaultfontfeatures{Ligatures=TeX}
\setmainfont{PlusJakartaSans-Medium.ttf}[Path=../../../fonts/,BoldFont=PlusJakartaSans-Bold.ttf]
% Maths en sans-serif (Fira Math) avec chiffres et lettres droites de Plus
% Jakarta, pour que les formules se fondent dans le texte.
\usepackage[math-style=ISO,bold-style=ISO]{unicode-math}
\setmathfont{FiraMath-Regular.otf}[Path=../../../fonts/]
\setmathfont{PlusJakartaSans-Medium.ttf}[Path=../../../fonts/,range={up/num,up/{Latin,latin},\mathup}]
\usepackage{icomma} % virgule décimale sans espace en mode math
% Caractères mathématiques Unicode absents des polices de texte (Plus Jakarta,
% Archivo Black) : rendus par leur équivalent mathématique, en texte comme en
% formule — sinon ils s'affichent en carré vide (ex. « Arithmétique dans ℤ »).
\usepackage{newunicodechar}
\newunicodechar{ℝ}{\ensuremath{\mathbb{R}}}
\newunicodechar{ℤ}{\ensuremath{\mathbb{Z}}}
\newunicodechar{ℂ}{\ensuremath{\mathbb{C}}}
\newunicodechar{ℕ}{\ensuremath{\mathbb{N}}}
\newunicodechar{ℚ}{\ensuremath{\mathbb{Q}}}
\newunicodechar{𝔻}{\ensuremath{\mathbb{D}}}
\newunicodechar{α}{\ensuremath{\alpha}}
\newunicodechar{β}{\ensuremath{\beta}}
\newunicodechar{γ}{\ensuremath{\gamma}}
\newunicodechar{δ}{\ensuremath{\delta}}
\newunicodechar{ε}{\ensuremath{\varepsilon}}
\newunicodechar{θ}{\ensuremath{\theta}}
\newunicodechar{λ}{\ensuremath{\lambda}}
\newunicodechar{μ}{\ensuremath{\mu}}
\newunicodechar{σ}{\ensuremath{\sigma}}
\newunicodechar{τ}{\ensuremath{\tau}}
\newunicodechar{φ}{\ensuremath{\varphi}}
\newunicodechar{ψ}{\ensuremath{\psi}}
\newunicodechar{ω}{\ensuremath{\omega}}
\newunicodechar{Σ}{\ensuremath{\Sigma}}
% majuscules produites par \MakeUppercase dans les titres (α-aminés -> Α)
\newunicodechar{Α}{\ensuremath{\alpha}}
\newunicodechar{Β}{\ensuremath{\beta}}
\newunicodechar{Γ}{\ensuremath{\Gamma}}
\newunicodechar{Δ}{\ensuremath{\Delta}}
\newunicodechar{Ε}{\ensuremath{\varepsilon}}
\newunicodechar{Θ}{\ensuremath{\Theta}}
\newunicodechar{Λ}{\ensuremath{\Lambda}}
\newunicodechar{Μ}{\ensuremath{\mu}}
\newunicodechar{Τ}{\ensuremath{\tau}}
\newunicodechar{Φ}{\ensuremath{\Phi}}
\newunicodechar{Ψ}{\ensuremath{\Psi}}
\newunicodechar{Ω}{\ensuremath{\Omega}}
\newunicodechar{↦}{\ensuremath{\mapsto}}
\newunicodechar{∈}{\ensuremath{\in}}
\newunicodechar{∉}{\ensuremath{\notin}}
\newunicodechar{∧}{\ensuremath{\wedge}}
\newunicodechar{⇔}{\ensuremath{\Leftrightarrow}}
\newunicodechar{⟺}{\ensuremath{\Longleftrightarrow}}
\newunicodechar{⇒}{\ensuremath{\Rightarrow}}
\newunicodechar{⟹}{\ensuremath{\Longrightarrow}}
\newunicodechar{∘}{\ensuremath{\circ}}
\newunicodechar{∪}{\ensuremath{\cup}}
\newunicodechar{∩}{\ensuremath{\cap}}
\newunicodechar{≡}{\ensuremath{\equiv}}
\newunicodechar{⊂}{\ensuremath{\subset}}
\newunicodechar{⊥}{\ensuremath{\perp}}
\newunicodechar{∃}{\ensuremath{\exists}}
\newunicodechar{∀}{\ensuremath{\forall}}
\newunicodechar{∛}{\ensuremath{\sqrt[3]{\,}}}
\newunicodechar{∜}{\ensuremath{\sqrt[4]{\,}}}
\newunicodechar{⃗}{\ensuremath{{}^{\to}}}
\newunicodechar{ⁿ}{\ifmmode^{n}\else\textsuperscript{\ensuremath{n}}\fi}
\newunicodechar{ˣ}{\ifmmode^{x}\else\textsuperscript{\ensuremath{x}}\fi}
\newunicodechar{ᵇ}{\ifmmode^{b}\else\textsuperscript{\ensuremath{b}}\fi}
\newunicodechar{ʳ}{\ifmmode^{r}\else\textsuperscript{\ensuremath{r}}\fi}
\newunicodechar{ᵘ}{\ifmmode^{u}\else\textsuperscript{\ensuremath{u}}\fi}
\newunicodechar{ʸ}{\ifmmode^{y}\else\textsuperscript{\ensuremath{y}}\fi}
\newunicodechar{ᵃ}{\ifmmode^{a}\else\textsuperscript{\ensuremath{a}}\fi}
\newunicodechar{ᵉ}{\ifmmode^{e}\else\textsuperscript{\ensuremath{e}}\fi}
\newunicodechar{ᵏ}{\ifmmode^{k}\else\textsuperscript{\ensuremath{k}}\fi}
\newunicodechar{ᵖ}{\ifmmode^{p}\else\textsuperscript{\ensuremath{p}}\fi}
\newunicodechar{ᴺ}{\ifmmode^{N}\else\textsuperscript{\ensuremath{N}}\fi}
\newunicodechar{ᶜ}{\ifmmode^{c}\else\textsuperscript{\ensuremath{c}}\fi}
\newunicodechar{ʰ}{\ifmmode^{h}\else\textsuperscript{\ensuremath{h}}\fi}
\newunicodechar{ᵠ}{\ifmmode^{q}\else\textsuperscript{\ensuremath{q}}\fi}
\newunicodechar{ₙ}{\ifmmode_{n}\else\textsubscript{\ensuremath{n}}\fi}
\newunicodechar{ₐ}{\ifmmode_{a}\else\textsubscript{\ensuremath{a}}\fi}
\newunicodechar{ᵢ}{\ifmmode_{i}\else\textsubscript{\ensuremath{i}}\fi}
\newunicodechar{ᵣ}{\ifmmode_{r}\else\textsubscript{\ensuremath{r}}\fi}
\newunicodechar{ₖ}{\ifmmode_{k}\else\textsubscript{\ensuremath{k}}\fi}
\newunicodechar{ᵦ}{\ifmmode_{\beta}\else\textsubscript{\ensuremath{\beta}}\fi}
\newunicodechar{ₓ}{\ifmmode_{x}\else\textsubscript{\ensuremath{x}}\fi}
\newfontfamily\popdisplay{ArchivoBlack-Regular.ttf}[Path=../../../fonts/]
\newfontfamily\poplogo{SpaceGrotesk-Bold.ttf}[Path=../../../fonts/]
\newfontfamily\popsemi{PlusJakartaSans-SemiBold.ttf}[Path=../../../fonts/]
\newfontfamily\popxbold{PlusJakartaSans-ExtraBold.ttf}[Path=../../../fonts/]
\newfontfamily\popxboldls{PlusJakartaSans-ExtraBold.ttf}[Path=../../../fonts/,LetterSpace=3.0]

\setlist[enumerate]{leftmargin=1.7em,itemsep=5pt,topsep=4pt}
\setlist[itemize]{leftmargin=1.7em,itemsep=4pt,topsep=4pt,label={\color{poporange}\textbullet}}
\setlength{\parindent}{0pt}
\setlength{\parskip}{5pt}
\renewcommand{\arraystretch}{1.25}

% pgfplots : style Pop par défaut
\pgfplotsset{
  every axis/.append style={axis line style={popink,line width=1pt},tick style={popink},
    grid style={popline,line width=0.5pt},label style={font=\small},tick label style={font=\footnotesize}},
  popcurve/.style={poporange,line width=1.8pt,no markers,smooth},
}
% Même style disponible dans un \draw TikZ simple (hors pgfplots)
\tikzset{popcurve/.style={poporange,line width=1.8pt}}
\tikzset{popgrid/.style={popline,very thin}}
\colorlet{popcurve}{poporange} % le nom du style est parfois utilisé comme couleur

% ============================================================
% Pill / squiggle
% ============================================================
\newcommand{\poppill}[3]{%
  \tikz[baseline=-0.55ex]\node[draw=popink,line width=1.2pt,rounded corners=2.6pt,%
    fill=#2,inner xsep=6.5pt,inner ysep=3pt]{%
    \popxboldls\fontsize{7}{7}\selectfont\color{#3}\MakeUppercase{#1}};%
}
\newcommand{\popsquiggle}[1]{%
  \tikz[baseline=-0.4ex]{
    \draw[#1,line width=2.6pt,line cap=round]
      plot [smooth] coordinates {(0,0.09) (0.34,0.01) (0.68,0.21) (1.02,0.01) (1.36,0.10)};
  }%
}
% Mot-clé surligné (marqueur orange sous le texte)
\newcommand{\cle}[1]{{\popxbold\color{popdark}#1}}

% ============================================================
% En-tête / pied
% ============================================================
\pagestyle{fancy}
\fancyhf{}
\renewcommand{\headrulewidth}{0pt}
\renewcommand{\footrulewidth}{0pt}
\lhead{\raisebox{-0.15ex}{\poplogo\fontsize{13}{13}\selectfont\color{popink}OnBuch\color{poporange}+}}
\rhead{\poppill{\DOCMATIERE\ \textbullet\ \DOCNIVEAU\ \textbullet\ Leçon \DOCLECON}{white}{popink}}
\lfoot{\popxbold\fontsize{7.6}{7.6}\selectfont\color{popmuted}\MakeUppercase{Cours signé OnBuch+}\ \textbullet\ {\popsemi\DOCTITRE}}
\rfoot{\tikz[baseline=-0.6ex]\node[rounded corners=8.5pt,fill=popink,minimum height=17pt,minimum width=17pt,inner xsep=5pt,inner sep=0]{\popxbold\fontsize{8}{8}\selectfont\color{popcream}\thepage\,/\,\pageref*{LastPage}};}

% ============================================================
% Titres
% ============================================================
\newcommand{\chaptitre}[2]{%
  \begin{tcolorbox}[enhanced,colback=poporange,colframe=popink,boxrule=2.6pt,arc=11pt,
      drop shadow=popink,left=15pt,right=15pt,top=12pt,bottom=11pt,before skip=3pt,after skip=18pt]
    \poppill{#2}{popink}{popcream}\par\vspace{9pt}
    {\raggedright\hyphenpenalty=10000\exhyphenpenalty=10000\popdisplay\fontsize{22}{24}\selectfont\color{popcream}\MakeUppercase{#1}\par}
  \end{tcolorbox}
}
% Section numérotée : pastille orange avec le numéro + titre Archivo
\newcounter{popsec}
\newcommand{\coursec}[1]{%
  \stepcounter{popsec}\setcounter{popsub}{0}%
  \par\vspace{16pt}\noindent
  \begin{minipage}{\linewidth}
  \tikz[baseline=-0.7ex]\node[draw=popink,line width=1.6pt,fill=poporange,rounded corners=4pt,minimum size=22pt,inner sep=0]{\popdisplay\fontsize{12}{12}\selectfont\color{popcream}\thepopsec};%
  \hspace{8pt}{\popdisplay\fontsize{15}{17}\selectfont\color{popink}\MakeUppercase{#1}}\par
  \vspace{4pt}\noindent{\color{poporange}\rule{36pt}{3.6pt}}
  \end{minipage}\par\nopagebreak\vspace{8pt}
}
\newcounter{popsub}
\newcommand{\courssub}[1]{%
  \stepcounter{popsub}%
  \par\vspace{9pt}\noindent{\popxbold\fontsize{11.5}{13}\selectfont\color{popink}{\color{poporange}\thepopsec.\thepopsub}\hspace{6pt}#1}\par\nopagebreak\vspace{4pt}
}

% ============================================================
% Boîtes pédagogiques — même recette : fond blanc, bordure épaisse colorée,
% ombre dure encre, badge plein. Titre optionnel [#1].
% ============================================================
\tcbset{popbase/.style={breakable,enhanced,colback=white,boxrule=2.3pt,arc=9pt,
  shadow={3pt}{-3pt}{0pt}{popink},left=14pt,right=14pt,top=11pt,bottom=11pt,before skip=11pt,after skip=15pt,
  fontupper=\popsemi\fontsize{10.6}{14.5}\selectfont\color{popink},
  fontlower=\popsemi\fontsize{10.6}{14.5}\selectfont\color{popink}}}
% Coche dessinée (le glyphe ✓ n'existe pas dans les polices du document)
\newcommand{\popcheck}[1]{\tikz[baseline=-0.5ex]\draw[#1,line width=1.6pt,line cap=round,line join=round](-0.13,0.0)--(-0.04,-0.09)--(0.13,0.1);}
\providecommand{\coche}{\popcheck{popgreen}}
\providecommand{\resultat}[1]{\par\smallskip\noindent\fcolorbox{popgreen}{popgreenL}{\parbox{\dimexpr\linewidth-2\fboxsep-2\fboxrule\relax}{\textbf{Résultat.} #1}}\par\smallskip}
\newcommand{\popboxhead}[3]{\poppill{#1}{#2}{white}\ifx\relax#3\relax\else\hspace{6pt}{\popxbold\fontsize{10.6}{12}\selectfont\color{popink}#3}\fi\par\vspace{7pt}}

\newtcolorbox{aretenir}[1][]{popbase,colframe=popgreen,before upper={\popboxhead{À retenir}{popgreen}{#1}}}
\newtcolorbox{exemplebox}[1][]{popbase,colframe=poporange,before upper={\popboxhead{Exemple}{poporange}{#1}}}
\newtcolorbox{definition}[1][]{popbase,colframe=popblue,before upper={\popboxhead{Définition}{popblue}{#1}}}
\newtcolorbox{propriete}[1][]{popbase,colframe=poppurple,before upper={\popboxhead{Propriété}{poppurple}{#1}}}
\newtcolorbox{methode}[1][]{popbase,colframe=popgold,before upper={\popboxhead{Méthode}{popgold}{#1}}}
\newtcolorbox{attention}[1][]{popbase,colframe=poppink,before upper={\popboxhead{Attention}{poppink}{#1}}}
\newtcolorbox{experience}[1][]{popbase,colframe=popblue,colback=popblueL,before upper={\popboxhead{Expérience}{popblue}{#1}}}
% « Le savais-tu ? » : carte violette pleine (contexte, culture, Cameroun)
\newtcolorbox{savaistu}[1][]{popbase,colback=poppurple,colframe=popink,
  fontupper=\popsemi\fontsize{10.4}{14.2}\selectfont\color{white},
  before upper={\poppill{Le savais-tu\,?}{white}{poppurple}\ifx\relax#1\relax\else\hspace{6pt}{\popxbold\color{white}#1}\fi\par\vspace{7pt}}}
% Exercice résolu : énoncé en haut, \tcblower puis la solution sur fond crème
\newcounter{popexo}\newcounter{popexr}
\newtcolorbox{exoresolu}[1][]{popbase,colframe=poporange,colbacklower=popcream,
  segmentation style={popink,line width=1.2pt,dashed},
  before upper={\stepcounter{popexr}\popboxhead{Exercice résolu \thepopexr}{poporange}{#1}},
  before lower={\poppill{Solution}{popink}{popcream}\par\vspace{6pt}}}
% Exercice (non résolu en place) — la correction est dans la partie Corrigés
\newtcolorbox{exercice}[1][]{popbase,colframe=popink,
  before upper={\stepcounter{popexo}\popboxhead{Exercice \thepopexo}{popink}{#1}}}
% Corrigé
\newtcolorbox{corrige}[1][]{popbase,colframe=popgreen,colback=popgreenL,
  before upper={\popboxhead{Corrigé}{popgreen}{#1}}}
% Objectifs / prérequis / bilan
\newtcolorbox{objectifs}{popbase,colframe=popink,colback=poporangeL,
  before upper={\popboxhead{Objectifs de la leçon}{popink}{}}}
\newtcolorbox{prerequis}{popbase,colframe=popink,colback=popblueL,
  before upper={\popboxhead{Prérequis}{popblue}{}}}
\newtcolorbox{bilan}{popbase,colframe=popink,colback=white,boxrule=2.8pt,
  before upper={\popboxhead{Fiche bilan}{poporange}{L'essentiel à connaître pour l'examen}}}
% Figure encadrée avec légende
\newtcolorbox{popfigure}[1][]{enhanced,breakable=false,colback=white,colframe=popline,boxrule=1.4pt,arc=8pt,
  left=10pt,right=10pt,top=10pt,bottom=8pt,before skip=10pt,after skip=12pt,halign=center,
  lower separated=false,halign lower=center,
  fontlower=\popsemi\fontsize{9}{11.5}\selectfont\color{popmuted},
  #1}
\newcommand{\legende}[1]{\par\vspace{6pt}{\popsemi\fontsize{9}{11.5}\selectfont\color{popmuted}#1}}

% ============================================================
% Signature OnBuch+
% ============================================================
\newcommand{\poplogoinline}[1]{{\poplogo\fontsize{#1}{#1}\selectfont\color{popink}OnBuch\color{poporange}+}}
\newcommand{\signatureonbuch}{%
  \begin{tcolorbox}[enhanced,colback=popink,colframe=popink,boxrule=2.2pt,arc=10pt,
      drop shadow=poporange,left=14pt,right=14pt,top=10pt,bottom=10pt,before skip=0pt,after skip=0pt]
    \tikz[baseline=-0.7ex]\node[circle,fill=popgreen,draw=popcream,line width=1.3pt,minimum size=22pt,inner sep=0]{\popcheck{white}};%
    \hspace{9pt}\begin{minipage}[c]{0.62\linewidth}
      {\popxbold\fontsize{12}{13}\selectfont\color{popcream}Signé\ \textbullet\ {\poplogo OnBuch\color{poporange}+}}\par\vspace{2pt}
      {\raggedright\popsemi\fontsize{8}{10.5}\selectfont\color{popline}\MakeUppercase{Équipe pédagogique OnBuch+}\\\MakeUppercase{Conforme au programme officiel MINESEC}\par}
    \end{minipage}\hfill
    {\poplogo\fontsize{22}{22}\selectfont\color{popcream}OnBuch\color{poporange}+}
  \end{tcolorbox}%
}

% ============================================================
% Page de garde
% #1 titre · #2 matière · #3 niveau · #4 module · #5 n° leçon · #6 durée ·
% #7 description / objectifs (texte ou itemize)
% ============================================================
\newcommand{\pagegarde}[7]{%
  \thispagestyle{empty}
  \newgeometry{a4paper,margin=1.7cm,top=1.6cm,bottom=1.4cm}%
  \begin{tikzpicture}[remember picture,overlay]
    \fill[poporange!18] ([xshift=-3.2cm,yshift=-3.6cm]current page.north east) circle (4.6cm);
    \fill[poppurple!14] ([xshift=2.4cm,yshift=7.6cm]current page.south west) circle (3.8cm);
  \end{tikzpicture}%
  {\poplogo\fontsize{30}{30}\selectfont\color{popink}OnBuch\color{poporange}+}\hfill
  \raisebox{0.6ex}{\poppill{Cours signé \textbullet\ Premium}{popink}{popcream}}\par
  \vspace{1pt}\popsquiggle{poppurple}\par
  \vspace{8pt}
  \begin{tcolorbox}[enhanced,colback=poporange,colframe=popink,boxrule=2.8pt,arc=13pt,
      drop shadow=popink,left=18pt,right=18pt,top=12pt,bottom=13pt,after skip=10pt]
    \poppill{#2 \textbullet\ #3}{popink}{popcream}\hspace{5pt}\poppill{#4}{white}{popink}\par\vspace{8pt}
    {\popdisplay\fontsize{14}{14}\selectfont\color{popink}LEÇON #5}\par\vspace{4pt}
    {\raggedright\hyphenpenalty=10000\exhyphenpenalty=10000\popdisplay\fontsize{25}{27}\selectfont\color{popcream}\MakeUppercase{#1}\par}
    \vspace{8pt}
    {\popxbold\fontsize{9.5}{9.5}\selectfont\color{popink}\MakeUppercase{Durée conseillée : #6}}
  \end{tcolorbox}
  \begin{tcolorbox}[enhanced,colback=white,colframe=popink,boxrule=2.2pt,arc=10pt,
      drop shadow=popink,left=16pt,right=16pt,top=10pt,bottom=10pt,after skip=8pt]
    \poppill{Dans cette leçon}{poppurple}{white}\par\vspace{5pt}
    \popsemi\fontsize{9.3}{12.6}\selectfont\color{popink}#7
  \end{tcolorbox}
  \noindent\begin{minipage}[t]{0.487\linewidth}
    \begin{tcolorbox}[enhanced,colback=popgreen,colframe=popink,boxrule=2.2pt,arc=10pt,
        drop shadow=popink,left=11pt,right=11pt,top=10pt,bottom=10pt,height=3.1cm,valign=center]
      \tcbox[colback=white,colframe=popink,boxrule=1.3pt,arc=5pt,boxsep=0pt,nobeforeafter,
        left=3pt,right=3pt,top=3pt,bottom=3pt,valign=center]{\qrcode[height=1.9cm]{https://play.google.com/store/apps/details?id=com.luvvix.onbuch&hl=fr}}%
      \hspace{8pt}\begin{minipage}[c]{0.5\linewidth}
        {\popdisplay\fontsize{11}{12.5}\selectfont\color{white}\MakeUppercase{Télécharge OnBuch}}\par\vspace{4pt}
        {\raggedright\popsemi\fontsize{8.4}{11}\selectfont\color{white}Gratuit \textbullet\ Google Play\\Tuteur IA, annales, concours, résultats\par}
      \end{minipage}
    \end{tcolorbox}
  \end{minipage}\hfill
  \begin{minipage}[t]{0.487\linewidth}
    \begin{tcolorbox}[enhanced,colback=poppurple,colframe=popink,boxrule=2.2pt,arc=10pt,
        drop shadow=popink,left=11pt,right=11pt,top=10pt,bottom=10pt,height=3.1cm,valign=center]
      \tikz[baseline=-0.7ex]\node[circle,fill=white,draw=popink,line width=1.3pt,minimum size=1.6cm,inner sep=0]{\popdisplay\fontsize{24}{24}\selectfont\color{poppurple}+};%
      \hspace{8pt}\begin{minipage}[c]{0.6\linewidth}
        {\popdisplay\fontsize{11}{12.5}\selectfont\color{white}\MakeUppercase{Passe à OnBuch+}}\par\vspace{4pt}
        {\raggedright\popsemi\fontsize{8.4}{11}\selectfont\color{white}Tous les cours signés, Tuteur IA illimité, fiches PDF — onglet OnBuch+ de l'app\par}
      \end{minipage}
    \end{tcolorbox}
  \end{minipage}
  \vspace{8pt}
  \popcontenu
  \vfill
  \signatureonbuch
  \restoregeometry
  \newpage
}

% Rangée « Ce cours contient » (page de garde)
\newcommand{\poptile}[3]{% #1 couleur #2 gros texte #3 légende
  \begin{tcolorbox}[enhanced,colback=#1,colframe=popink,boxrule=2pt,arc=9pt,shadow={2.5pt}{-2.5pt}{0pt}{popink},
      left=6pt,right=6pt,top=9pt,bottom=9pt,halign=center,valign=center,height=1.75cm,width=\linewidth,nobeforeafter]
    {\popdisplay\fontsize{12.5}{13}\selectfont\color{white}\MakeUppercase{#2}}\par\vspace{3pt}
    {\centering\popsemi\fontsize{7.8}{9.5}\selectfont\color{white}#3\par}
  \end{tcolorbox}}
\newcommand{\popcontenu}{%
  \par\noindent{\popxboldls\fontsize{7.5}{7.5}\selectfont\color{popmuted}CE COURS SIGNÉ CONTIENT}\par\vspace{7pt}
  \noindent\begin{minipage}[t]{0.235\linewidth}\poptile{popblue}{Cours}{complet, illustré, pas à pas}\end{minipage}\hfill
  \begin{minipage}[t]{0.235\linewidth}\poptile{poporange}{Méthodes}{et exercices résolus}\end{minipage}\hfill
  \begin{minipage}[t]{0.235\linewidth}\poptile{poppink}{Exercices}{type examen + corrigés}\end{minipage}\hfill
  \begin{minipage}[t]{0.235\linewidth}\poptile{popgreen}{Fiche bilan}{l'essentiel à retenir}\end{minipage}\par}

% Sommaire
\newtcolorbox{sommaire}{enhanced,colback=white,colframe=popink,boxrule=2.2pt,arc=10pt,
  drop shadow=popink,left=18pt,right=18pt,top=13pt,bottom=13pt,before skip=6pt,after skip=20pt,
  before upper={\poppill{Sommaire}{poppurple}{white}\par\vspace{9pt}\popsemi\fontsize{10.6}{14.5}\selectfont\color{popink}}}

% Signature de fin de cours — poussée tout en bas de la dernière page
\newcommand{\findecours}{%
  \par\vfill\nopagebreak
  \begin{center}\popsquiggle{poporange}\end{center}\vspace{4pt}
  \signatureonbuch
}
\AtBeginDocument{\def\llbracket{\mathopen{[\mkern-3.5mu[}}\def\rrbracket{\mathclose{]\mkern-3.5mu]}}} % intervalles d'entiers (glyphe ⟦ absent de la police)
% Glyphes absents de Fira Math : redéfinis à partir de glyphes disponibles
\AtBeginDocument{%
  \def\setminus{\mathbin{\backslash}}\let\smallsetminus\setminus
  \def\vdots{\mathord{\vbox{\baselineskip3.2pt\lineskiplimit0pt\kern4pt\hbox{.}\hbox{.}\hbox{.}}}}%
  \def\ddots{\mathinner{\mkern1mu\raise6.5pt\hbox{.}\mkern2mu\raise3.6pt\hbox{.}\mkern2mu\raise0.7pt\hbox{.}\mkern1mu}}%
  \def\triangle{\mathord{\scalebox{0.9}{$\symup{\Delta}$}}}%
  \def\nabla{\mathord{\raisebox{\depth}{\scalebox{1}[-1]{$\symup{\Delta}$}}}}%
  \def\checkmark{\popcheck{popdark}}%
  \def\star{\mathbin{\ast}}\def\bigstar{\mathord{\ast}}%
  \def\diamond{\mathbin{\scalebox{0.6}{$\Diamond$}}}%
  \def\bigcup{\mathop{\mathchoice{\raisebox{-0.3em}{\scalebox{1.7}{$\cup$}}}{\scalebox{1.25}{$\cup$}}{\cup}{\cup}}\displaylimits}%
  \def\bigcap{\mathop{\mathchoice{\raisebox{-0.3em}{\scalebox{1.7}{$\cap$}}}{\scalebox{1.25}{$\cap$}}{\cap}{\cap}}\displaylimits}%
  \def\lesseqqgtr{\mathrel{\lessgtr}}\def\bigoplus{\mathop{\oplus}\displaylimits}%
}
\providecommand{\ssi}{si et seulement si} % abréviation fréquente
\AtBeginDocument{\providecommand{\wideparen}{\overparen}} % arc de cercle

\begin{document}

\pagegarde{Classification des éléments}{PCT}{3ème}{Module 1 : Chimie}{N02}{1 séance (1 h chacune)}{Cette leçon présente le tableau périodique de Mendeleïev, outil fondamental de classification des éléments chimiques. L'élève apprend à y repérer les lignes (périodes), les colonnes (familles) et les éléments usuels de la vie courante camerounaise.
\vspace{4pt}
\begin{multicols}{2}\small
\begin{itemize}
\item À la fin de cette leçon, je sais définir un élément chimique et rappeler son symbole.
\item À la fin de cette leçon, je sais décrire l'organisation générale du tableau de Mendeleïev (lignes et colonnes).
\item À la fin de cette leçon, je sais définir les notions de période et de famille (groupe).
\item À la fin de cette leçon, je sais repérer un élément dans le tableau à partir de son symbole ou de son numéro atomique.
\item À la fin de cette leçon, je sais citer quelques familles remarquables (alcalins, halogènes, gaz nobles) et leurs propriétés communes.
\item À la fin de cette leçon, je sais identifier les éléments usuels (H, O, C, N, Fe, Al, Cu, Ca, Na, Cl...) et leur présence dans la vie courante.
\end{itemize}\end{multicols}}

\begin{sommaire}
\begin{multicols}{2}
\begin{enumerate}
\item Rappels : éléments chimiques et symboles
\item Historique et construction du tableau de Mendeleïev
\item Structure du tableau : lignes (périodes) et colonnes (familles)
\item Quelques familles remarquables
\item Éléments usuels et applications dans la vie courante
\item Activité d'intégration
\item Méthodes et exercices résolus
\item Exercices
\item Corrigés des exercices
\item Fiche bilan
\end{enumerate}
\end{multicols}
\end{sommaire}

\begin{objectifs}
\begin{itemize}
\item À la fin de cette leçon, je sais définir un élément chimique et rappeler son symbole.
\item À la fin de cette leçon, je sais décrire l'organisation générale du tableau de Mendeleïev (lignes et colonnes).
\item À la fin de cette leçon, je sais définir les notions de période et de famille (groupe).
\item À la fin de cette leçon, je sais repérer un élément dans le tableau à partir de son symbole ou de son numéro atomique.
\item À la fin de cette leçon, je sais citer quelques familles remarquables (alcalins, halogènes, gaz nobles) et leurs propriétés communes.
\item À la fin de cette leçon, je sais identifier les éléments usuels (H, O, C, N, Fe, Al, Cu, Ca, Na, Cl...) et leur présence dans la vie courante.
\item À la fin de cette leçon, je sais lire les informations portées dans une case du tableau (symbole, numéro atomique Z, masse molaire).
\item À la fin de cette leçon, je sais appliquer les consignes de sécurité lors de la manipulation de produits contenant certains éléments (eau de Javel, acides).
\end{itemize}
\end{objectifs}

\begin{prerequis}
\begin{itemize}
\item Notion d'atome et de molécule (leçon précédente de 3ème)
\item Symboles chimiques des éléments courants (H, O, C, Fe...)
\item Notion de corps pur simple et de corps pur composé
\item Lecture d'un tableau à double entrée (compétence mathématique)
\end{itemize}
\end{prerequis}

\newpage

\coursec{Rappels : éléments chimiques et symboles}

Au marché Mokolo à Yaoundé, le vendeur de quincaillerie range les clous, les vis et les écrous dans des casiers séparés pour s'y retrouver rapidement. De la même façon, les chimistes du XIXe siècle se sont retrouvés face à plus de 60 éléments connus sans aucun ordre logique. En 1869, le Russe Dmitri Mendeleïev a eu l'idée géniale de ranger les éléments comme un jeu de cartes, en laissant même des cases vides pour des éléments pas encore découverts. Comment ce tableau est-il organisé, et comment l'utiliser pour retrouver les informations sur les éléments qui nous entourent, comme l'aluminium des marmites ou le fer des tiges à béton ?

\courssub{Qu'est-ce qu'un élément chimique ?}

Tout ce qui nous entoure — l'air que nous respirons, l'eau que nous buvons, la marmite en aluminium, le câble en cuivre, le fer à béton de nos maisons — est constitué de \textbf{matériaux}. En chimie, on distingue les \emph{corps purs} (une seule substance) des \emph{mélanges} (plusieurs substances). Un \textbf{élément chimique} est un corps pur fondamental qui ne peut pas être décomposé en substances plus simples par des moyens chimiques ordinaires.

\begin{definition}[Élément chimique]
Un \textbf{élément chimique} est un ensemble d'atomes identiques caractérisés par le même \textbf{numéro atomique} $Z$ (nombre de protons dans le noyau). Il existe 118 éléments connus à ce jour, dont une grande majorité sont des métaux.
\end{definition}

\courssub{Le symbole chimique : une écriture universelle}

Pour communiquer sans ambiguïté dans toutes les langues, chaque élément possède un \textbf{symbole chimique} normalisé par l'IUPAC (Union internationale de chimie pure et appliquée).

\begin{definition}[Symbole chimique]
Le \textbf{symbole chimique} d'un élément s'écrit avec \textbf{une ou deux lettres} :
\begin{itemize}
    \item La \textbf{première lettre est toujours une majuscule}.
    \item La \textbf{deuxième lettre, si elle existe, est toujours une minuscule}.
\end{itemize}
Exemples : \ce{H} (hydrogène), \ce{O} (oxygène), \ce{Fe} (fer), \ce{Na} (sodium), \ce{Cl} (chlore).
\end{definition}

\begin{attention}[Piège majeur : Majuscule / Minuscule]
Ne confonds jamais la majuscule et la minuscule dans un symbole à deux lettres !
\begin{itemize}
    \item \ce{Co} = \textbf{Cobalt} (métal magnétique, utilisé dans les alliages).
    \item \ce{CO} = \textbf{Monoxyde de carbone} (gaz toxique, \emph{molécule} formée de Carbone \ce{C} et d'Oxygène \ce{O}).
\end{itemize}
Écrire « \ce{CO} » pour le cobalt est une \textbf{erreur grave} qui change complètement le sens chimique. Au BEPC, cette confusion est sanctionnée.
\end{attention}

\courssub{Origine des symboles : français vs latin}

La plupart des symboles dérivent du nom français ou anglais (première lettre du nom). Mais pour des éléments connus depuis l'Antiquité, le symbole vient du \textbf{latin} (ou du grec). C'est une source fréquente d'erreurs : il faut les apprendre par cœur.

\begin{exemplebox}[Symboles d'origine latine]
\begin{center}
\begin{tabular}{|l|l|l|}\hline
\rowcolor{popblueL}
\textbf{Élément (français)} & \textbf{Nom latin} & \textbf{Symbole} \\ \hline
Fer & Ferrum & \ce{Fe} \\
Cuivre & Cuprum & \ce{Cu} \\
Argent & Argentum & \ce{Ag} \\
Or & Aurum & \ce{Au} \\
Plomb & Plumbum & \ce{Pb} \\
Étain & Stannum & \ce{Sn} \\
Sodium & Natrium & \ce{Na} \\
Potassium & Kalium & \ce{K} \\
Mercure & Hydrargyrum & \ce{Hg} \\ \hline
\end{tabular}
\end{center}
\emph{Astuce mnémotechnique :} « \textbf{Fe}rrum = \textbf{Fe}r », « \textbf{Na}trium = \textbf{Na} (soude) », « \textbf{K}alium = \textbf{K} (potasse) ».
\end{exemplebox}

\courssub{Les 15 éléments usuels à connaître par cœur}

Au BEPC, tu dois identifier instantanément le nom et le symbole des éléments les plus courants dans ton environnement. Voici la liste de référence pour la classe de 3ème.

\begin{popfigure}
\begin{center}
\begin{tabularx}{\linewidth}{|l|X|l|X|}\hline
\rowcolor{popblueL}
\textbf{Élément} & \textbf{Symbole} & \textbf{Élément} & \textbf{Symbole} \\ \hline
Hydrogène & \ce{H} & Soufre & \ce{S} \\
Carbone & \ce{C} & Chlore & \ce{Cl} \\
Azote & \ce{N} & Potassium & \ce{K} \\
Oxygène & \ce{O} & Calcium & \ce{Ca} \\
Sodium & \ce{Na} & Fer & \ce{Fe} \\
Magnésium & \ce{Mg} & Cuivre & \ce{Cu} \\
Aluminium & \ce{Al} & Zinc & \ce{Zn} \\ \hline
\end{tabularx}
\end{center}
\legende{Tableau des 15 éléments usuels à maîtriser (nom / symbole).}
\end{popfigure}

\begin{methode}[Apprendre les symboles efficacement]
\begin{enumerate}
    \item \textbf{Groupe par origine :} Sépare ceux qui viennent du français (\ce{H, C, N, O, S, Cl, Al, Zn}) de ceux qui viennent du latin (\ce{Na, K, Fe, Cu, Ca, Mg} -- attention \ce{Mg} vient de Magnésia mais symbole moderne).
    \item \textbf{Cartes mentales :} Associe chaque symbole à un objet du quotidien (voir ci-dessous).
    \item \textbf{Teste-toi :} Cache la colonne « Symbole » et écris-les sur une ardoise. Répète jusqu'à 0 faute.
\end{enumerate}
\end{methode}

\courssub{Éléments du quotidien au Cameroun}

La chimie n'est pas abstraite : elle est dans ta cuisine, ton atelier, ton corps.

\begin{savaistu}[Composition du corps humain]
Le corps humain est une « usine chimique » vivante. En masse, il est constitué d'environ :
\begin{itemize}
    \item \textbf{65 \% d'Oxygène (\ce{O})} (majoritairement dans l'eau \ce{H2O}).
    \item \textbf{18 \% de Carbone (\ce{C})} (base de toutes les molécules organiques : protéines, sucres, graisses, ADN).
    \item \textbf{10 \% d'Hydrogène (\ce{H})} (dans l'eau et les molécules organiques).
    \item \textbf{3 \% d'Azote (\ce{N})} (protéines, ADN).
    \item Les \textbf{4 \% restants} : Calcium (\ce{Ca}) pour les os et dents, Phosphore (\ce{P}), Potassium (\ce{K}), Soufre (\ce{S}), Sodium (\ce{Na}), Chlore (\ce{Cl}), Magnésium (\ce{Mg}), Fer (\ce{Fe}) \dots
\end{itemize}
Sans ces éléments, pas de vie ! L'anémie, fréquente chez les jeunes au Cameroun, est souvent un manque de \textbf{Fer (\ce{Fe})} (viande rouge, feuilles de manioc, haricots).
\end{savaistu}

\begin{center}
\begin{tabularx}{\linewidth}{|l|l|X|}\hline
\rowcolor{popblueL}
\textbf{Élément} & \textbf{Objet / Source au Cameroun} & \textbf{Rôle / Utilisation} \\ \hline
\ce{Al} (Aluminium) & Marmites, casseroles, papier alu, toits en tôle & Léger, bon conducteur thermique, ne rouille pas (couche d'alumine protectrice). \\
\ce{Fe} (Fer) & Tiges à béton (fer à béton), portails, motos, outils & Résistance mécanique, magnétique. Rouille (\ce{Fe2O3}) au contact air+eau $\to$ peinture/galvanisation. \\
\ce{Cu} (Cuivre) & Câbles électriques ENEO, fils moteurs, tuyaux plomberie & Meilleur conducteur électrique après l'argent, malléable, antibactérien. \\
\ce{Zn} (Zinc) & Galvanisation (toitures, seaux), piles (Zn/C), crème solaire & Protège le fer de la rouille (anode sacrificielle). \\
\ce{Ca} (Calcium) & Lait, fromage, os, coquilles d'œufs, chaux (blanchiment cases) & Solidité osseuse, transmission nerveuse, chaux vive \ce{CaO} pour mortier. \\
\ce{Na} (Sodium) & Sel de table (\ce{NaCl}), savon (soude \ce{NaOH}) & Goût, conservation aliments, saponification (fabrication savon). \\
\ce{K} (Potassium) & Engrais (NPK), bananes, ignames, cacao & Croissance plantes (engrais), contraction musculaire/cardiaque. \\ \hline
\end{tabularx}
\end{center}

\courssub{Exercice résolu : Identifier éléments et symboles}

\begin{exoresolu}[Reconnaissance d'éléments et pièges]
\textbf{Énoncé :} Parmi les notations suivantes, lesquelles représentent un \textbf{élément chimique} (symbole correct) ? Pour celles qui ne sont pas des symboles d'éléments, explique l'erreur.
\[
\ce{H}, \quad \ce{co}, \quad \ce{CO}, \quad \ce{Fe}, \quad \ce{CL}, \quad \ce{Na}, \quad \ce{Al}, \quad \ce{ca}
\]

\tcblower
\textbf{Solution détaillée :}

\begin{enumerate}
    \item \ce{H} : \textbf{Correct}. Hydrogène (1 lettre, majuscule).
    \item \ce{co} : \textbf{Incorrect}. Deux minuscules. Le symbole du Cobalt s'écrit \ce{Co} (C majuscule, o minuscule).
    \item \ce{CO} : \textbf{Incorrect pour un élément}. Deux majuscules. Cela représente la \textbf{molécule} de monoxyde de carbone (Carbone + Oxygène). Piège classique !
    \item \ce{Fe} : \textbf{Correct}. Fer (symbole latin \emph{Ferrum}). Majuscule + minuscule.
    \item \ce{CL} : \textbf{Incorrect}. Deuxième lettre en majuscule. Le Chlore s'écrit \ce{Cl} (C majuscule, l minuscule).
    \item \ce{Na} : \textbf{Correct}. Sodium (symbole latin \emph{Natrium}). Majuscule + minuscule.
    \item \ce{Al} : \textbf{Correct}. Aluminium. Majuscule + minuscule.
    \item \ce{ca} : \textbf{Incorrect}. Deux minuscules. Le Calcium s'écrit \ce{Ca} (C majuscule, a minuscule).
\end{enumerate}

\textbf{Bilan :} Seules \ce{H}, \ce{Fe}, \ce{Na}, \ce{Al} sont des symboles d'éléments correctement écrits.
\end{exoresolu}

\begin{exercice}[Entraînement rapide]
\begin{enumerate}
    \item Donne le symbole de : Magnésium, Soufre, Potassium, Cuivre, Zinc.
    \item Donne le nom de l'élément correspondant à : \ce{Ca}, \ce{K}, \ce{Fe}, \ce{Cl}, \ce{Mg}.
    \item Corrige les erreurs d'écriture : \ce{fe}, \ce{NA}, \ce{al}, \ce{CL}, \ce{cu}.
    \item Un électricien à Douala utilise du fil de \ce{Cu} et du fil d'\ce{Al}. Pourquoi le cuivre est-il préféré pour les installations intérieures malgré son prix plus élevé ?
\end{enumerate}
\end{exercice}

\begin{corrige}[Exercice n°1]
\begin{enumerate}
    \item \ce{Mg}, \ce{S}, \ce{K}, \ce{Cu}, \ce{Zn}.
    \item Calcium, Potassium, Fer, Chlore, Magnésium.
    \item \ce{Fe}, \ce{Na}, \ce{Al}, \ce{Cl}, \ce{Cu} (respecter majuscule/minuscule !).
    \item Le cuivre (\ce{Cu}) a une \textbf{conductivité électrique} supérieure à l'aluminium (\ce{Al}) (environ 1,6 fois meilleure). Il chauffe moins pour le même courant, supporte mieux les courts-circuits, et ses connexions sont plus fiables (moins d'oxydation superficielle). L'aluminium est utilisé pour les lignes aériennes haute tension (ENEO) car il est \textbf{plus léger} et \textbf{moins cher} pour de grandes portées.
\end{enumerate}
\end{corrige}

\begin{aretenir}[Ce qu'il faut retenir de cette section]
\begin{itemize}
    \item Un \textbf{élément chimique} = ensemble d'atomes de même numéro atomique $Z$.
    \item \textbf{Symbole} = 1 majuscule \textbf{ou} 1 majuscule + 1 minuscule (jamais 2 majuscules, jamais 2 minuscules).
    \item \textbf{15 symboles usuels} à connaître par cœur (dont \ce{Fe, Cu, Na, K, Ca, Zn, Al, Mg} d'origine latine ou spécifique).
    \item \textbf{Contexte camerounais :} \ce{Al} (marmites), \ce{Fe} (fer à béton), \ce{Cu} (câbles ENEO), \ce{Ca} (os/lait), \ce{Na} (sel), \ce{K} (engrais/banane).
    \item \textbf{Piège BEPC :} \ce{Co} (Cobalt) $\neq$ \ce{CO} (Monoxyde de carbone).
\end{itemize}
\end{aretenir}

\coursec{Historique et construction du tableau de Mendeleïev}

Dans la section précédente, nous avons appris qu'un \cle{élément chimique} est un ensemble d'atomes identiques, caractérisé par son \cle{numéro atomique} $Z$, et que chaque élément possède un \cle{symbole} : \ce{H} pour l'hydrogène, \ce{O} pour l'oxygène, \ce{Fe} pour le fer, \ce{Au} pour l'or. Mais voici une question que se posaient les savants il y a 150 ans : comment \textbf{ranger} tous ces éléments ? Y a-t-il un ordre caché dans la nature, comme les notes d'une gamme de musique ? C'est à cette question qu'un chimiste russe, Dmitri Mendeleïev, a apporté une réponse géniale en 1869. Suivons son raisonnement pas à pas.

\courssub{Le problème historique : un désordre de plus de 60 éléments}

\courssub{Une situation de la vie courante : la bibliothèque sans classement}

Imagine que tu entres dans une grande bibliothèque de Yaoundé où les livres sont posés au hasard : un roman entre un dictionnaire et un livre de cuisine, un atlas sous une pile de bandes dessinées. Retrouver un livre serait un cauchemar ! Le bibliothécaire résout le problème en \textbf{classant} : les livres de sciences ensemble, les romans ensemble, et dans chaque rayon, un ordre précis.

Au milieu du XIX\textsuperscript{e} siècle, les chimistes vivaient exactement cette situation avec les éléments chimiques.

\begin{experience}[Document historique : la situation de la chimie vers 1860]
\textbf{Observation (document) :} vers 1860, on connaît plus de \textbf{60 éléments} : l'or, l'argent, le cuivre et le fer (connus depuis l'Antiquité), mais aussi des éléments découverts récemment comme l'aluminium (1825) ou le sodium (1807). Chaque chimiste les range à sa manière : certains par couleur, d'autres par densité, d'autres encore par « famille » de métaux.

\textbf{Problème :} aucun classement ne rend compte de \emph{toutes} les propriétés à la fois. On ne peut ni prévoir les propriétés d'un élément à partir de sa place, ni deviner s'il reste des éléments à découvrir.

\textbf{Question scientifique :} existe-t-il un critère unique permettant de ranger tous les éléments de façon logique ?
\end{experience}

Plusieurs savants avaient déjà tenté leur chance : Döbereiner avait remarqué des « triades » (groupes de trois éléments aux propriétés voisines, comme chlore, brome et iode), et Newlands avait proposé une « loi des octaves » (les propriétés se répéteraient tous les huit éléments, comme les notes de musique). Mais aucune de ces idées ne fonctionnait pour \emph{tous} les éléments connus. Il fallait un génie pour voir plus loin.

\courssub{l'idée géniale de Dmitri Mendeleïev}

\courssub{Le classement par masse atomique croissante}

Dmitri Ivanovitch Mendeleïev (1834-1907), professeur de chimie à l'université de Saint-Pétersbourg en Russie, préparait un manuel de chimie pour ses étudiants. Pour organiser son livre, il écrivit sur des cartes le nom de chaque élément, sa \textbf{masse atomique} et ses propriétés principales (aspect, densité, réaction avec l'oxygène, formule de son oxyde...). Puis il rangea ces cartes par \textbf{masse atomique croissante}.

\begin{experience}[La démarche de Mendeleïev, reconstituée]
\textbf{Hypothèse :} la masse atomique pourrait être le critère de classement.

\textbf{Expérience de pensée :} Mendeleïev aligne les éléments de la plus légère à la plus lourde : \ce{H} (1), \ce{Li} (7), \ce{Be} (9), \ce{B} (11), \ce{C} (12), \ce{N} (14), \ce{O} (16), \ce{F} (19), \ce{Na} (23), \ce{Mg} (24), \ce{Al} (27), \ce{Si} (28), \ce{P} (31), \ce{S} (32), \ce{Cl} (35,5), \ce{K} (39)...

\textbf{Observation décisive :} les propriétés \textbf{se répètent régulièrement} ! Le lithium \ce{Li}, le sodium \ce{Na} et le potassium \ce{K} sont tous des métaux mous qui réagissent violemment avec l'eau. Le fluor \ce{F} et le chlore \ce{Cl} sont tous deux des gaz très réactifs. C'est comme si, en avançant dans la liste, on retrouvait périodiquement « le même caractère ».

\textbf{Conclusion :} en disposant les éléments en lignes et en plaçant les uns sous les autres ceux qui se ressemblent, on obtient un tableau où chaque \textbf{colonne} rassemble des éléments aux propriétés semblables.
\end{experience}

\begin{definition}[Classification périodique]
La \cle{classification périodique} est un rangement de tous les éléments chimiques dans un tableau, tel que :
\begin{itemize}
\item les éléments sont ordonnés selon un critère croissant (masse atomique chez Mendeleïev, numéro atomique $Z$ aujourd'hui) ;
\item les éléments placés dans une même \textbf{colonne} possèdent des \textbf{propriétés chimiques semblables}.
\end{itemize}
On dit que les propriétés des éléments sont une \emph{fonction périodique} de leur rang : elles se répètent à intervalles réguliers.
\end{definition}

\begin{savaistu}[Une idée venue en jouant aux cartes ?]
Selon la légende, Mendeleïev était un passionné du jeu de \emph{patience} (la « réussite »), un jeu de cartes solitaire où l'on doit ranger les cartes par couleur et par valeur. Il aurait fabriqué un « jeu de cartes des éléments » : une carte par élément, avec sa masse atomique et ses propriétés. En jouant à les ranger sur sa table, comme une patience, il aurait vu apparaître l'ordre caché des éléments. Une autre version raconte que le tableau lui serait apparu \textbf{en rêve}, après des jours de travail acharné. Quelle que soit la vérité, une chose est sûre : en février-mars 1869, il présenta sa classification à la Société chimique de Russie sous le titre « Essai d'un système des éléments d'après leurs poids atomiques et leurs fonctions chimiques ».
\end{savaistu}

\courssub{Les cases vides : une prédiction audacieuse}

Voici le geste le plus audacieux de Mendeleïev. En rangeant ses cartes, il remarqua que certains éléments ne « tombaient » pas à la bonne place : leurs propriétés ne correspondaient pas à celles de leur colonne. Plutôt que de déformer son tableau, il fit un pari extraordinaire :

\begin{propriete}[Le pari des cases vides]
Mendeleïev laissa des \textbf{cases vides} dans son tableau et affirma que ces cases correspondaient à des éléments \textbf{encore inconnus}, qui seraient découverts un jour. Mieux : grâce à la place de chaque case vide, il \textbf{prédit les propriétés} de ces éléments manquants (masse atomique, densité, aspect, réactions chimiques), par comparaison avec leurs voisins de ligne et de colonne.
\end{propriete}

Il baptisa ces éléments hypothétiques avec le préfixe sanskrit \emph{eka-} (« un »), qui signifie « celui qui est juste en dessous » :
\begin{itemize}
\item \textbf{eka-aluminium} : l'élément qui devait se trouver \emph{sous} l'aluminium \ce{Al} ;
\item \textbf{eka-silicium} : l'élément qui devait se trouver \emph{sous} le silicium \ce{Si}.
\end{itemize}

\begin{attention}[Prédire n'est pas deviner]
La prédiction de Mendeleïev n'était pas de la voyance : c'était un raisonnement scientifique. Si les propriétés varient régulièrement le long d'une ligne et d'une colonne, alors un élément situé entre deux voisins connus doit avoir des propriétés \emph{intermédiaires}. C'est le même raisonnement que tu utilises quand tu complètes une suite de nombres : si une case se trouve entre 5 et 7, tu prédis 6.
\end{attention}

\courssub{La vérification expérimentale : le gallium et le germanium}

En science, une théorie n'est acceptée que si ses \textbf{prédictions sont confirmées par l'expérience}. C'est exactement ce qui arriva à Mendeleïev, de son vivant, à deux reprises.

\courssub{la découverte du gallium}

En 1875, le chimiste français Paul-Émile Lecoq de Boisbaudran découvrit un nouvel élément dans un minerai de zinc, et le baptisa \textbf{gallium} (\ce{Ga}), en l'honneur de la Gaule. Quand Mendeleïev lut la description de ce métal, il écrivit immédiatement : c'est l'eka-aluminium ! Comparons les prédictions de 1869 et les mesures de 1875.

\begin{center}
\rowcolors{1}{popcream}{white}
\begin{tabularx}{\linewidth}{|l|X|X|}
\hline
\rowcolor{popblueL} \textbf{Propriété} & \textbf{Eka-aluminium (prédit en 1869)} & \textbf{Gallium (mesuré en 1875)} \\
\hline
Masse atomique & environ 68 & 69,7 \\
\hline
Densité & 5,9 & 5,94 \\
\hline
Aspect & métal, point de fusion bas & métal qui fond à \SI{29,8}{\celsius} (dans la main !) \\
\hline
Formule de l'oxyde & \ce{Ea2O3} & \ce{Ga2O3} \\
\hline
Réaction avec les acides & se dissout lentement & se dissout lentement \\
\hline
\end{tabularx}
\end{center}

\begin{exemplebox}[Une prédiction presque parfaite]
Mendeleïev avait prédit une \textbf{densité de 5,9} pour l'eka-aluminium. Le gallium découvert en 1875 a une densité mesurée de \textbf{5,94}. L'écart est de seulement $5,94 - 5,9 = 0,04$, soit moins de 1 \% d'erreur ! Mendeleïev alla même plus loin : la première mesure de Lecoq de Boisbaudran donnait une densité de 4,7, et Mendeleïev lui écrivit qu'il s'était trompé. Le Français refit sa mesure avec un échantillon plus pur... et trouva 5,94. Le théoricien russe avait raison contre l'expérimentateur français !
\end{exemplebox}

\courssub{la découverte du germanium}

En 1886, le chimiste allemand Clemens Winkler découvrit le \textbf{germanium} (\ce{Ge}). Là encore, les propriétés mesurées correspondaient remarquablement à celles prédites pour l'eka-silicium : masse atomique prédite 72 contre 72,6 mesurée ; densité prédite 5,5 contre 5,35 mesurée. La classification de Mendeleïev était définitivement validée : elle ne servait pas seulement à \emph{ranger} les éléments, elle permettait de \emph{prévoir}.

\begin{popfigure}
\begin{tikzpicture}[scale=1, every node/.style={font=\small}]
% Ligne de temps
\draw[line width=2pt, popblue, -{Stealth[length=3mm]}] (0,0) -- (14.6,0);
% Graduations et événements
\foreach \x/\annee/\txt/\col in {
  1.2/1869/{Mendeleïev publie sa\\classification et prédit\\des éléments inconnus}/popblue,
  4.9/1875/{Découverte du gallium \ce{Ga}\\(Lecoq de Boisbaudran)\\confirme l'eka-aluminium}/popgreen,
  8.6/1886/{Découverte du germanium \ce{Ge}\\(Winkler) confirme\\l'eka-silicium}/poporange,
  12.6/{Aujourd'hui}/{Tableau moderne :\\118 éléments classés\\par numéro atomique $Z$}/poppurple}{
  \draw[line width=1.5pt, \col] (\x,-0.18) -- (\x,0.18);
  \node[above=0.15cm, align=center, text=\col, font=\bfseries] at (\x,0.18) {\annee};
  \node[below=0.25cm, align=center, text=\col] at (\x,-0.18) {\txt};
}
\end{tikzpicture}
\legende{Frise chronologique : de la classification de Mendeleïev (1869) au tableau moderne. Les découvertes du gallium (1875) et du germanium (1886) ont confirmé les prédictions et rendu la classification célèbre dans le monde entier.}
\end{popfigure}

\courssub{Du tableau de Mendeleïev au tableau moderne}

Le tableau de Mendeleïev comportait quelques imperfections : certains éléments semblaient « inversés » quand on suivait strictement l'ordre des masses atomiques (par exemple, le tellure \ce{Te} et l'iode \ce{I} : pour que l'iode soit avec le chlore et le brome, il fallait le placer \emph{avant} le tellure, plus léger). Mendeleïev pensait que les masses étaient mal mesurées. En réalité, c'est le critère de classement lui-même qui devait changer.

Au début du XX\textsuperscript{e} siècle, les travaux de Rutherford puis de Moseley (1913) montrèrent que chaque atome est caractérisé par la \textbf{charge de son noyau}, c'est-à-dire son nombre de protons : le \cle{numéro atomique} $Z$.

\begin{aretenir}[Le critère du tableau moderne]
Dans le tableau périodique \textbf{moderne}, les éléments sont rangés par \cle{numéro atomique} $Z$ \textbf{croissant} (et non plus par masse atomique). Avec ce critère, toutes les « inversions » disparaissent : l'iode ($Z = 53$) vient naturellement après le tellure ($Z = 52$) et se retrouve dans la famille des halogènes, avec le fluor et le chlore.
\end{aretenir}

\begin{attention}[Piège classique au BEPC]
Ne confonds pas les deux critères de classement :
\begin{itemize}
\item \textbf{Mendeleïev (1869)} : classement par \textbf{masse atomique} croissante ;
\item \textbf{Tableau moderne} : classement par \textbf{numéro atomique} $Z$ croissant.
\end{itemize}
Dans la plupart des cas, les deux ordres coïncident (car la masse augmente en général avec $Z$), mais pas toujours. Si un exercice demande « selon quel critère les éléments sont-ils rangés dans le tableau actuel ? », la bonne réponse est : \emph{par numéro atomique $Z$ croissant}.
\end{attention}

\begin{propriete}[Le tableau périodique actuel]
Le tableau périodique actuel compte \textbf{118 éléments}, numérotés de $Z = 1$ (hydrogène \ce{H}) à $Z = 118$ (oganesson \ce{Og}). Parmi eux, \textbf{94 éléments sont naturels} (présents sur Terre, dans les roches, l'eau, l'air ou les êtres vivants) ; les autres ont été fabriqués artificiellement en laboratoire et sont pour la plupart très instables.
\end{propriete}

\begin{exoresolu}[Critère de classement et prédiction]
\textbf{Énoncé.} 1) Quel critère Mendeleïev a-t-il utilisé en 1869 pour ranger les éléments ? 2) Quel critère est utilisé dans le tableau moderne ? 3) Mendeleïev avait prédit pour l'eka-silicium une densité de 5,5. Le germanium découvert en 1886 a une densité de 5,35. Calculer l'écart entre la prédiction et la mesure, puis conclure.
\tcblower
\textbf{Solution détaillée.}
\begin{enumerate}
\item En 1869, Mendeleïev a rangé les éléments par \textbf{masse atomique croissante}, en plaçant dans une même colonne les éléments aux propriétés chimiques semblables.
\item Dans le tableau moderne, les éléments sont rangés par \textbf{numéro atomique $Z$ croissant}.
\item L'écart vaut : $5,5 - 5,35 = 0,15$. Cet écart est très faible (moins de 3 \% de la valeur mesurée). \textbf{Conclusion :} la prédiction de Mendeleïev était remarquablement exacte, ce qui confirme que sa classification traduit un ordre réel de la nature.
\end{enumerate}
\end{exoresolu}

\begin{exercice}[À toi de jouer]
1) Pourquoi Mendeleïev a-t-il laissé des cases vides dans son tableau ? 2) Que signifie le nom « eka-aluminium » ? 3) Citer les deux découvertes qui ont confirmé les prédictions de Mendeleïev, avec leurs dates. 4) Le tableau moderne compte 118 éléments. Combien sont naturels ?
\end{exercice}

\begin{corrige}[Exercice : à toi de jouer]
1) Parce que, pour respecter la régularité des propriétés, certaines places devaient rester libres : Mendeleïev a parié qu'elles correspondraient à des éléments encore inconnus, dont il a prédit les propriétés. 2) « Eka-aluminium » signifie « l'élément situé juste en dessous de l'aluminium » (\emph{eka} = « un » en sanskrit). 3) La découverte du gallium en 1875 (qui confirme l'eka-aluminium) et celle du germanium en 1886 (qui confirme l'eka-silicium). 4) 94 éléments sont naturels ; les 24 autres sont artificiels.
\end{corrige}

\begin{aretenir}[L'essentiel de la section]
\begin{itemize}
\item Vers 1860, plus de 60 éléments étaient connus, sans classement logique.
\item En \textbf{1869}, Mendeleïev range les éléments par \textbf{masse atomique croissante} et constate que les propriétés se répètent \textbf{périodiquement}.
\item Il laisse des \textbf{cases vides} et prédit les propriétés d'éléments inconnus (eka-aluminium, eka-silicium).
\item La découverte du \textbf{gallium} (1875, densité 5,94 contre 5,9 prédite) et du \textbf{germanium} (1886) confirme ses prédictions.
\item Le tableau \textbf{moderne} classe les éléments par \textbf{numéro atomique $Z$ croissant} ; il compte \textbf{118 éléments}, dont \textbf{94 naturels}.
\end{itemize}
\end{aretenir}

\coursec{Structure du tableau : lignes (périodes) et colonnes (familles)}

Nous avons vu dans la section précédente comment Mendeleïev, en rangeant les éléments par masse croissante et en regroupant ceux qui se ressemblent, a construit son célèbre tableau. Il nous reste maintenant à apprendre à \emph{lire} ce tableau, comme on apprend à lire une carte géographique : il faut connaître ses lignes, ses colonnes, ses zones et ses cases. C'est exactement l'objectif de cette section. À la fin, tu sauras localiser n'importe quel élément, comme tu saurais localiser un quartier de Yaoundé sur un plan de la ville.

\courssub{Vue d'ensemble : une carte à deux entrées}

Observe le tableau périodique affiché dans la salle de sciences de ton collège. Ce que tu vois ressemble à un grand casier, comme celui où l'on range les bouteilles de boissons au marché : des cases alignées en lignes horizontales et en colonnes verticales. Chaque case contient un élément chimique, et la place de chaque case n'est pas due au hasard : elle obéit à des règles précises.

\begin{definition}[Période et famille]
Dans le tableau de classification périodique :
\begin{itemize}
\item une \cle{période} est une \textbf{ligne horizontale} du tableau. Le tableau moderne compte \textbf{7 périodes}, numérotées de 1 à 7, de haut en bas ;
\item une \cle{famille} (ou \cle{groupe}) est une \textbf{colonne verticale} du tableau. Le tableau moderne compte \textbf{18 colonnes}, numérotées de 1 à 18, de gauche à droite.
\end{itemize}
\end{definition}

\begin{aretenir}[L'essentiel de la structure]
Le tableau périodique est organisé en \textbf{7 lignes horizontales} (les périodes) et \textbf{18 colonnes verticales} (les familles ou groupes). Chaque élément occupe une case située à l'intersection d'une période et d'une colonne.
\end{aretenir}

\courssub{Les périodes : les lignes horizontales}

\textbf{Observation et interprétation}\par

Prenons trois éléments que tu connais déjà : le sodium \ce{Na} (présent dans le sel de cuisine), le magnésium \ce{Mg} et le chlore \ce{Cl} (utilisé pour désinfecter l'eau de boisson). Dans le tableau, ces trois éléments sont rangés sur la \emph{même ligne} : la troisième. Or, si l'on dessine le modèle de Bohr de leurs atomes, on constate qu'ils possèdent tous \textbf{trois couches électroniques} : K, L et M.

Ce n'est pas une coïncidence. C'est la règle de construction des lignes du tableau.

\begin{propriete}[Signification du numéro de période]
Le numéro de la \cle{période} d'un élément est égal au \textbf{nombre de couches électroniques occupées} dans l'atome de cet élément.
\begin{itemize}
\item Période 1 : les atomes ont 1 couche électronique (couche K) ;
\item Période 2 : les atomes ont 2 couches électroniques (K et L) ;
\item Période 3 : les atomes ont 3 couches électroniques (K, L et M) ; et ainsi de suite jusqu'à la période 7.
\end{itemize}
\end{propriete}

\begin{exemplebox}[Lire les périodes]
\begin{itemize}
\item L'hydrogène \ce{H} ($Z = 1$) et l'hélium \ce{He} ($Z = 2$) sont seuls sur la période 1 : leurs atomes n'ont que la couche K.
\item Le carbone \ce{C} ($Z = 6$) est sur la période 2 : son atome possède les couches K et L.
\item Le fer \ce{Fe} ($Z = 26$) est sur la période 4 : son atome possède quatre couches électroniques.
\end{itemize}
\end{exemplebox}

\begin{attention}[La période 1 est très courte]
Ne t'étonne pas de voir seulement deux cases sur la première ligne : la période 1 ne contient que l'hydrogène et l'hélium. Les périodes suivantes sont plus longues. Cela vient du nombre d'électrons que chaque couche peut accueillir.
\end{attention}

\courssub{Les familles : les colonnes verticales}

\textbf{Pourquoi ranger en colonnes ?}\par

Revenons à l'idée de Mendeleïev : il a empilé les éléments qui \emph{se comportent de la même façon}. Par exemple, le lithium \ce{Li}, le sodium \ce{Na} et le potassium \ce{K} sont trois métaux mous, légers, qui réagissent violemment avec l'eau. Au laboratoire, on les conserve d'ailleurs dans l'huile pour éviter tout contact avec l'humidité de l'air. Mendeleïev les a placés les uns sous les autres, dans la même colonne.

\begin{propriete}[Signification d'une famille]
Les éléments d'une même \cle{famille} (même colonne) possèdent des \textbf{propriétés chimiques voisines} : ils réagissent de manière semblable avec les autres substances. Cette ressemblance s'explique par le fait que leurs atomes possèdent le \textbf{même nombre d'électrons sur leur couche électronique externe}.
\end{propriete}

\begin{exemplebox}[Deux familles à connaître]
\begin{itemize}
\item \textbf{Colonne 1} : la famille des \cle{métaux alcalins} (lithium \ce{Li}, sodium \ce{Na}, potassium \ce{K}...). Ce sont des métaux mous qui réagissent violemment avec l'eau.
\item \textbf{Colonne 18} : la famille des \cle{gaz nobles} (hélium \ce{He}, néon \ce{Ne}, argon \ce{Ar}...). Ce sont des gaz très peu réactifs ; on utilise l'argon dans certaines ampoules et le néon dans les enseignes lumineuses.
\end{itemize}
Nous étudierons ces familles en détail dans la section suivante.
\end{exemplebox}

\courssub{Schéma d'ensemble du tableau}

La figure ci-dessous représente de façon simplifiée l'organisation du tableau : les 7 périodes, les 18 colonnes, et la grande frontière qui sépare les métaux des non-métaux.

\begin{popfigure}
\begin{center}
\begin{tikzpicture}[scale=0.62, every node/.style={transform shape}]
% Grille des cases : 18 colonnes x 7 périodes, avec cases vides en haut
% Période 1 : H (col1) et He (col18)
\fill[popgoldL] (0,6) rectangle (1,7);
\node at (0.5,6.5) {\scriptsize \ce{H}};
\fill[popgreenL] (17,6) rectangle (18,7);
\node at (17.5,6.5) {\scriptsize \ce{He}};
% Période 2 : col1-2 métaux, col13-18 non-métaux
\fill[popblueL] (0,5) rectangle (1,6);
\fill[popblueL] (1,5) rectangle (2,6);
\foreach \c in {12,13,14,15,16,17} {\fill[popgreenL] (\c,5) rectangle (\c+1,6);}
% Période 3
\fill[popblueL] (0,4) rectangle (1,5);
\fill[popblueL] (1,4) rectangle (2,5);
\foreach \c in {12,13,14,15,16,17} {\fill[popgreenL] (\c,4) rectangle (\c+1,5);}
% Périodes 4 à 7 : pleines
\foreach \r in {0,1,2,3} {
  \foreach \c in {0,...,17} {\fill[popblueL] (\c,\r) rectangle (\c+1,\r+1);}
}
% Zone non-métaux en haut à droite des périodes 4-7 (approximation escalier)
\fill[popgreenL] (16,3) rectangle (18,4);
\fill[popgreenL] (17,2) rectangle (18,3);
% Contours de la grille
\draw[popdark, thin] (0,0) grid (18,4);
\draw[popdark, thin] (0,4) grid (2,7);
\draw[popdark, thin] (12,4) grid (18,7);
% Frontière métaux / non-métaux (escalier épais)
\draw[poporange, very thick] (2,4) -- (12,4) -- (12,5) -- (13,5) -- (13,4) -- (14,4) -- (14,3) -- (15,3) -- (15,2) -- (16,2) -- (16,1) -- (17,1) -- (17,0);
% Numéros des périodes
\foreach \r/\p in {6/1,5/2,4/3,3/4,2/5,1/6,0/7} {
  \node[left, popdark] at (0,\r+0.5) {\scriptsize \p};
}
\node[rotate=90, popdark] at (-1.1,3.5) {\scriptsize Périodes};
% Numéros de colonnes (quelques-uns)
\foreach \c/\n in {0/1,1/2,10/11,16/17,17/18} {
  \node[above, popdark] at (\c+0.5,7) {\scriptsize \n};
}
\node[above, popdark] at (9,7.15) {\scriptsize Colonnes (familles)};
% Étiquettes de zones
\node[popblue, align=center] at (9,1.5) {\small \textbf{MÉTAUX}};
\node[popgreen!60!black, align=center] at (15.2,5.5) {\small \textbf{NON-MÉTAUX}};
% Flèche vers H
\draw[-{Stealth}, popgold!70!black, thick] (0.5,6.1) -- (0.5,5.3) node[below, align=center, text=popgold!70!black] {\scriptsize H placé\\ \scriptsize à part};
\end{tikzpicture}
\end{center}
\legende{Schéma simplifié du tableau périodique. En bleu : les métaux (la grande majorité des éléments, à gauche et au centre). En vert : les non-métaux (à droite). Le trait orange en escalier marque la frontière entre métaux et non-métaux. Les numéros 1 à 7 à gauche désignent les périodes (lignes) ; les numéros en haut désignent les colonnes (familles). L'hydrogène H occupe une position particulière, isolée en haut à gauche.}
\end{popfigure}

\courssub{Lecture d'une case du tableau}

Chaque case du tableau est comme la carte d'identité d'un élément. Elle contient trois informations essentielles qu'il faut savoir lire sans se tromper.

\begin{definition}[Contenu d'une case]
Chaque case du tableau périodique indique :
\begin{itemize}
\item le \cle{symbole} de l'élément, au centre (une ou deux lettres, ex. : \ce{Fe}) ;
\item le \cle{numéro atomique} $Z$, en haut : c'est le nombre de protons du noyau (et aussi le nombre d'électrons de l'atome) ;
\item la \cle{masse molaire atomique} $M$, en bas : c'est la masse d'une mole d'atomes de cet élément, exprimée en grammes par mole (\SI{}{g/mol}).
\end{itemize}
\end{definition}

\begin{popfigure}
\begin{center}
\begin{tikzpicture}[scale=1.6]
% Case
\draw[popdark, very thick, fill=popcream] (0,0) rectangle (3,3);
% Numéro atomique en haut à gauche
\node[anchor=north west, popblue] at (0.15,2.9) {\large \textbf{26}};
% Symbole au centre
\node[popdark] at (1.5,1.5) {\Huge \textbf{Fe}};
% Masse molaire en bas
\node[anchor=south, poporange!80!black] at (1.5,0.15) {\large 56};
% Légendes fléchées
\draw[-{Stealth}, popblue, thick] (0.4,2.75) -- (-1.6,3.6) node[left, align=right, text=popblue] {\textbf{Numéro atomique} $Z$\\ (nombre de protons)};
\draw[-{Stealth}, popdark, thick] (1.5,1.9) -- (4.6,2.6) node[right, align=left, text=popdark] {\textbf{Symbole} de l'élément\\ (ici le fer)};
\draw[-{Stealth}, poporange!80!black, thick] (1.5,0.35) -- (4.6,-0.5) node[right, align=left, text=poporange!80!black] {\textbf{Masse molaire atomique} $M$\\ en \SI{}{g/mol}};
\end{tikzpicture}
\end{center}
\legende{Zoom sur la case du fer \ce{Fe}. En haut : le numéro atomique $Z = 26$ (l'atome de fer possède 26 protons et 26 électrons). Au centre : le symbole Fe. En bas : la masse molaire atomique $M = \SI{56}{g/mol}$ (une mole d'atomes de fer pèse \SI{56}{g}).}
\end{popfigure}

\begin{attention}[Ne confonds pas Z et M !]
C'est l'erreur la plus fréquente au BEPC. Dans une case :
\begin{itemize}
\item le \textbf{numéro atomique} $Z$ est un \textbf{nombre entier, sans unité} : il compte les protons. Il est placé \textbf{en haut} ;
\item la \textbf{masse molaire} $M$ s'exprime en \textbf{g/mol} : c'est une masse. Elle est placée \textbf{en bas} et sa valeur est presque toujours plus grande que $Z$ (souvent environ le double).
\end{itemize}
Exemple : pour le fer, $Z = 26$ et $M = \SI{56}{g/mol}$. Écrire « le fer a 56 protons » est faux : le fer a 26 protons. Retiens : \emph{Z compte, M pèse}.
\end{attention}

\courssub{Repérer un élément dans le tableau}

Repérer un élément, c'est donner son « adresse » : sa période (ligne) et sa colonne (famille). C'est exactement comme donner l'adresse d'une maison : le quartier et la rue.

\begin{methode}[Repérer un élément à partir de son symbole ou de Z]
\begin{enumerate}
\item \textbf{Localise la case} de l'élément dans le tableau grâce à son symbole (ou à son numéro atomique $Z$, qui croît régulièrement de gauche à droite et de haut en bas).
\item \textbf{Lis le numéro de la ligne} (à gauche du tableau) : c'est le numéro de la \textbf{période}.
\item \textbf{Lis le numéro de la colonne} (en haut du tableau) : c'est le numéro de la \textbf{famille}.
\item \textbf{Conclus} en donnant les deux informations : « l'élément X se trouve dans la n\textsuperscript{e} période et la n\textsuperscript{e} colonne ».
\end{enumerate}
\end{methode}

\begin{exemplebox}[Le sodium et le chlore]
\begin{itemize}
\item Le sodium \ce{Na} ($Z = 11$) se trouve à l'intersection de la \textbf{3\textsuperscript{e} période} et de la \textbf{1\textsuperscript{re} colonne}. Son atome possède donc 3 couches électroniques.
\item Le chlore \ce{Cl} ($Z = 17$) se trouve à l'intersection de la \textbf{3\textsuperscript{e} période} et de la \textbf{17\textsuperscript{e} colonne}. Il est voisin du sodium sur la même ligne, mais leurs propriétés sont très différentes car ils ne sont pas dans la même famille : le sodium est un métal qui réagit avec l'eau, le chlore est un gaz non-métal désinfectant.
\end{itemize}
\end{exemplebox}

\begin{exoresolu}[Localiser l'aluminium]
Énoncé : l'aluminium \ce{Al}, utilisé pour fabriquer les marmites et les tôles, a pour numéro atomique $Z = 13$. À l'aide du tableau périodique, indique sa période et sa colonne, puis déduis-en le nombre de couches électroniques de son atome. L'aluminium est-il un métal ou un non-métal ?
\tcblower
Solution détaillée :
\begin{enumerate}
\item On repère la case d'abréviation \ce{Al} (numéro atomique 13) dans le tableau.
\item La case se trouve sur la \textbf{3\textsuperscript{e} ligne} : l'aluminium appartient à la \textbf{3\textsuperscript{e} période}.
\item La case se trouve dans la \textbf{13\textsuperscript{e} colonne} : l'aluminium appartient à la famille de la colonne 13.
\item Le numéro de période donne le nombre de couches électroniques : l'atome d'aluminium possède donc \textbf{3 couches électroniques} (K, L et M).
\item La case de l'aluminium est située \textbf{à gauche de la frontière en escalier} : l'aluminium est un \textbf{métal}. C'est cohérent avec la vie courante : il conduit la chaleur (marmites) et le courant électrique (câbles).
\end{enumerate}
Conclusion : \ce{Al} : 3\textsuperscript{e} période, 13\textsuperscript{e} colonne, 3 couches électroniques ; c'est un métal.
\end{exoresolu}

\courssub{Les grandes zones du tableau : métaux et non-métaux}

Un dernier coup d'œil sur la carte avant de conclure. Si tu colories le tableau, tu verrais qu'il se partage en deux grandes régions, séparées par une frontière en escalier.

\begin{definition}[Métaux et non-métaux dans le tableau]
\begin{itemize}
\item Les \cle{métaux} occupent la partie \textbf{gauche et centrale} du tableau : ce sont les éléments de \textbf{loin les plus nombreux} (environ les trois quarts des éléments). Exemples : fer \ce{Fe}, aluminium \ce{Al}, cuivre \ce{Cu}, zinc \ce{Zn}, sodium \ce{Na}, or \ce{Au}.
\item Les \cle{non-métaux} occupent la partie \textbf{droite} du tableau (au-dessus et à droite de la frontière en escalier). Exemples : carbone \ce{C}, azote \ce{N}, oxygène \ce{O}, chlore \ce{Cl}, soufre \ce{S}.
\end{itemize}
\end{definition}

\begin{savaistu}[Pourquoi l'hydrogène est-il placé à part ?]
L'hydrogène \ce{H} est un cas unique. On le dessine souvent en haut à gauche, au-dessus de la colonne 1, mais il \textbf{n'appartient à aucune famille} : ce n'est ni un métal alcalin, ni un gaz noble. C'est un gaz non-métal très léger, le plus léger de tous les éléments, et le plus abondant de l'Univers : c'est le « carburant » des étoiles comme notre Soleil. Pour marquer son originalité, de nombreux tableaux le placent dans une case isolée, à part des autres.
\end{savaistu}

\begin{attention}[Pièges de rédaction à éviter]
\begin{itemize}
\item Ne dis pas « la colonne 3 » quand tu veux parler de la 3\textsuperscript{e} ligne : une \textbf{ligne = période} (horizontale), une \textbf{colonne = famille} (verticale). Une astuce : les \emph{périodes} d'un match de football se succèdent dans le temps, comme les lignes se succèdent de haut en bas.
\item Deux éléments de la même \emph{ligne} n'ont pas forcément des propriétés voisines ; ce sont les éléments de la même \emph{colonne} qui se ressemblent chimiquement.
\item Le numéro de période donne le nombre de \textbf{couches} électroniques, pas le nombre d'électrons. Le sodium ($Z = 11$) a 11 électrons répartis sur 3 couches : il est en période 3.
\end{itemize}
\end{attention}

\begin{aretenir}[Bilan de la section]
\begin{itemize}
\item Le tableau périodique comporte \textbf{7 périodes} (lignes horizontales) et \textbf{18 familles} (colonnes verticales).
\item Le numéro de la période = nombre de \textbf{couches électroniques} de l'atome.
\item Les éléments d'une même colonne ont des \textbf{propriétés chimiques voisines}.
\item Chaque case indique le \textbf{symbole}, le \textbf{numéro atomique} $Z$ (en haut) et la \textbf{masse molaire} $M$ en \SI{}{g/mol} (en bas).
\item Les \textbf{métaux} (majorité) sont à gauche et au centre ; les \textbf{non-métaux} sont à droite ; l'\textbf{hydrogène} est placé à part.
\end{itemize}
\end{aretenir}

\begin{exercice}[Je m'entraîne]
\begin{enumerate}
\item Le potassium \ce{K} a pour numéro atomique $Z = 19$. À l'aide du tableau, donne sa période et sa colonne. À quelle famille appartient-il ? Prévois son comportement avec l'eau.
\item Le magnésium \ce{Mg} ($Z = 12$) et le béryllium \ce{Be} ($Z = 4$) sont dans la même colonne. Que peux-tu dire de leurs propriétés chimiques ? Combien de couches électroniques possède l'atome de magnésium ?
\item Dans la case du cuivre, on lit : 29 en haut, \ce{Cu} au centre, 63,5 en bas. Donne le nombre de protons, le nombre d'électrons et la masse molaire atomique du cuivre, en précisant les unités.
\end{enumerate}
\end{exercice}

\begin{corrige}[Exercice « Je m'entraîne »]
\begin{enumerate}
\item Le potassium \ce{K} ($Z = 19$) se trouve dans la \textbf{4\textsuperscript{e} période} et la \textbf{1\textsuperscript{re} colonne} : c'est un \textbf{métal alcalin}, comme le sodium. On prévoit donc qu'il réagit \textbf{violemment avec l'eau} (propriété commune à toute la famille).
\item \ce{Mg} et \ce{Be} étant dans la même colonne (colonne 2), ils ont des \textbf{propriétés chimiques voisines}. Le magnésium est dans la 3\textsuperscript{e} période : son atome possède \textbf{3 couches électroniques}.
\item Pour le cuivre : $Z = 29$, donc \textbf{29 protons} et \textbf{29 électrons} (l'atome est électriquement neutre) ; $M = \SI{63,5}{g/mol}$ : une mole d'atomes de cuivre pèse \SI{63,5}{g}.
\end{enumerate}
\end{corrige}

\coursec{Quelques familles remarquables}

Dans la section précédente, nous avons vu que le tableau de Mendeleïev est organisé en \cle{lignes} (les périodes) et en \cle{colonnes}. Mais pourquoi les chimistes ont-ils rangé les éléments en colonnes ? Parce qu'une colonne regroupe des éléments qui se ressemblent, un peu comme les membres d'une même famille partagent des traits communs. C'est cette idée que nous allons maintenant approfondir : certaines colonnes du tableau forment des \cle{familles chimiques} célèbres, dont tu croises les membres chaque jour sans le savoir — dans ta cuisine, à l'hôpital, dans les enseignes lumineuses de la ville.

\courssub{La notion de famille chimique}

\begin{definition}[Famille chimique]
Une \cle{famille chimique} est une colonne du tableau de classification périodique. Les éléments d'une même famille ont des \textbf{propriétés chimiques semblables} : ils réagissent de façon voisine avec les mêmes substances.
\end{definition}

Imagine une famille humaine : les frères et sœurs ne sont pas identiques, mais ils ont souvent le même caractère, les mêmes habitudes. De même, le sodium Na et le potassium K ne sont pas le même élément, mais tous deux réagissent violemment avec l'eau. Trois familles sont à connaître parfaitement pour le BEPC :

\begin{itemize}
\item la famille des \cle{alcalins} (1\iere{} colonne, sauf l'hydrogène H) : Li, Na, K\ldots{}
\item la famille des \cle{halogènes} (17\ieme{} colonne) : F, Cl, Br, I\ldots{}
\item la famille des \cle{gaz nobles} (18\ieme{} colonne) : He, Ne, Ar\ldots{}
\end{itemize}

\begin{popfigure}
\begin{center}
\begin{tikzpicture}[scale=0.62, every node/.style={font=\small}]
% Fond du tableau (18 colonnes, 7 lignes simplifiées)
\fill[popcream] (0,0) rectangle (18,7);
\draw[popline, thin] (0,0) grid (18,7);
% Colonne 1 : alcalins (sauf H)
\foreach \y/\e in {5/Li, 4/Na, 3/K, 2/Rb, 1/Cs}{
  \fill[poporangeL] (0,\y) rectangle (1,\y+1);
  \node at (0.5,\y+0.5) {\textbf{\e}};
}
\fill[white] (0,6) rectangle (1,7); \node at (0.5,6.5) {H};
% Colonne 17 : halogènes
\foreach \y/\e in {5/F, 4/Cl, 3/Br, 2/I}{
  \fill[popgreenL] (16,\y) rectangle (17,\y+1);
  \node at (16.5,\y+0.5) {\textbf{\e}};
}
% Colonne 18 : gaz nobles
\foreach \y/\e in {6/He, 5/Ne, 4/Ar, 3/Kr, 2/Xe}{
  \fill[popblueL] (17,\y) rectangle (18,\y+1);
  \node at (17.5,\y+0.5) {\textbf{\e}};
}
% Contours des familles
\draw[poporange, very thick] (0,1) rectangle (1,6);
\draw[popgreen, very thick] (16,2) rectangle (17,6);
\draw[popblue, very thick] (17,1) rectangle (18,7);
% Étiquettes
\node[poporange, font=\bfseries, align=center] at (0.5,0.3) {Alcalins};
\node[popgreen, font=\bfseries, align=center] at (16.5,1.3) {Halogènes};
\node[popblue, font=\bfseries, align=center] at (17.5,0.3) {Gaz nobles};
% Numéros de colonnes
\node[font=\footnotesize\itshape, popmuted] at (0.5,7.4) {1};
\node[font=\footnotesize\itshape, popmuted] at (16.5,7.4) {17};
\node[font=\footnotesize\itshape, popmuted] at (17.5,7.4) {18};
\end{tikzpicture}
\end{center}
\legende{Les trois familles remarquables dans le tableau périodique (représentation simplifiée) : en orange, les alcalins (1\iere{} colonne, sauf H) ; en vert, les halogènes (17\ieme{} colonne) ; en bleu, les gaz nobles (18\ieme{} colonne).}
\end{popfigure}

\begin{propriete}[Propriété des familles chimiques (admise)]
Les éléments d'une même colonne du tableau périodique ont des \textbf{propriétés chimiques voisines}. Cette propriété s'explique (tu le verras au lycée) par le fait que leurs atomes possèdent \textbf{le même nombre d'électrons sur leur couche externe}, et ce sont ces électrons externes qui participent aux réactions chimiques.
\end{propriete}

\begin{aretenir}
Retenir une famille, c'est retenir le « caractère » de tous ses membres d'un coup : si tu sais que le sodium réagit violemment avec l'eau, tu sais aussi que le potassium le fera. C'est toute la puissance de la classification de Mendeleïev.
\end{aretenir}

\courssub{La famille des alcalins : des métaux qui « explosent » au contact de l'eau}

\textbf{Présentation.} Les \cle{alcalins} occupent la première colonne du tableau, à l'exception de l'hydrogène H qui est un cas particulier. Ce sont le lithium Li, le sodium Na, le potassium K, le rubidium Rb et le césium Cs. Contrairement au fer ou au cuivre que tu connais, ce sont des métaux \textbf{mous} : on coupe le sodium au couteau, comme du \textbf{beurre} ! Ils brillent quand on vient de les couper, mais ce brillant disparaît vite au contact de l'air.

\textbf{Démarche d'investigation.} Pourquoi le professeur sort-il toujours le sodium d'un flacon rempli d'huile, et jamais d'un bocal d'eau ou d'une étagère à l'air libre ? Formulons une hypothèse : le sodium réagit avec l'eau, et peut-être même avec l'air. Vérifions par l'expérience.

\begin{experience}[Action du sodium sur l'eau (démonstration du professeur)]
\textbf{Matériel :} un cristallisoir rempli d'eau, un petit morceau de sodium conservé dans l'huile, une pince, du papier filtre, une allumette ou une baguette enflammée.\par
\textbf{Protocole :}
\begin{enumerate}
\item À l'aide d'une pince, prélever un petit morceau de sodium (taille d'un petit pois) et le sécher sur du papier filtre.
\item Le déposer délicatement à la surface de l'eau du cristallisoir.
\item Observer, puis approcher une flamme du gaz dégagé.
\end{enumerate}
\textbf{Observations :} le sodium fond en une petite bille argentée qui fuse à la surface de l'eau dans tous les sens, avec un sifflement ; il peut s'enflammer (flamme jaune-orangée). Le gaz dégagé produit une petite détonation à l'approche d'une flamme : c'est le \cle{dihydrogène} \ce{H2}. L'eau obtenue devient basique (elle contient de la soude \ce{NaOH}).\par
\textbf{Interprétation :} la réaction est vive et \textbf{exothermique} (elle dégage de la chaleur, ce qui fait fondre le sodium). Son équation s'écrit :
\[ \ce{2Na + 2H2O -> 2NaOH + H2} \]
\textbf{Conclusion :} le sodium réagit violemment avec l'eau ; c'est pourquoi on le \textbf{conserve dans l'huile}, qui l'isole de l'air humide et de l'eau.
\end{experience}

\begin{popfigure}
\begin{center}
\begin{tikzpicture}[scale=0.9, font=\small]
% Cristallisoir
\draw[popdark, thick] (-3,0) -- (-3,2.2);
\draw[popdark, thick] (3,0) -- (3,2.2);
\draw[popdark, thick] (-3,0) -- (3,0);
% Eau
\fill[popblueL, opacity=0.7] (-2.95,0.05) rectangle (2.95,1.6);
\draw[popblue] (-2.95,1.6) -- (2.95,1.6);
\node[popblue] at (2.2,0.5) {eau};
% Bille de sodium
\shade[ball color=popmuted] (0,1.75) circle (0.28);
\node[above, popdark] at (0,2.05) {bille de sodium};
% Trajectoire
\draw[poporange, dashed, -{Stealth}] (0.3,1.8) .. controls (1.2,2.1) .. (2.0,1.75);
% Flamme
\begin{scope}[shift={(2.0,1.75)}]
\fill[poporange] (0,0.15) .. controls (0.25,0.55) and (0.1,0.8) .. (0,1.0)
  .. controls (-0.1,0.8) and (-0.25,0.55) .. (0,0.15);
\fill[popgold] (0,0.2) .. controls (0.12,0.45) and (0.05,0.6) .. (0,0.72)
  .. controls (-0.05,0.6) and (-0.12,0.45) .. (0,0.2);
\node[poporange, below] at (0,-0.05) {flamme};
\end{scope}
% Bulles de H2
\foreach \p in {(-1.2,1.9), (-1.5,2.2), (-1.0,2.45), (-1.6,2.6)}{
  \draw[popblue] \p circle (0.09);
}
\node[popblue, align=center] at (-2.6,2.9) {dégagement\\ de \ce{H2}};
\draw[popblue, -{Stealth}] (-2.0,2.7) -- (-1.35,2.35);
% Étiquettes
\node[popdark, below] at (0,-0.25) {cristallisoir};
\end{tikzpicture}
\end{center}
\legende{Expérience du sodium sur l'eau : la bille de sodium fuse à la surface, dégage du dihydrogène \ce{H2} (bulles) et de la chaleur ; le gaz peut s'enflammer.}
\end{popfigure}

\begin{attention}[Sécurité : le sodium ne se manipule jamais à mains nues]
Le sodium réagit avec la \textbf{sueur} de la peau, qui contient de l'eau : il peut provoquer de graves brûlures. On le manipule toujours avec une \textbf{pince}, en petites quantités, et uniquement par le professeur. Au laboratoire, il est conservé \textbf{dans l'huile} (ou le pétrole), jamais dans l'eau ni à l'air libre.
\end{attention}

\begin{savaistu}[Du sel de cuisine au sodium métal]
Le sodium Na que contient ton sel de cuisine (\ce{NaCl}) n'a rien à voir avec le sodium métallique ! Dans le sel, le sodium est combiné au chlore sous forme d'ions et il est parfaitement inoffensif. C'est une leçon importante : les \textbf{propriétés d'un élément combiné} dans un corps composé ne sont pas celles du \textbf{corps simple}. Le potassium K, autre alcalin, est lui indispensable aux plantes : on le retrouve dans les engrais NPK utilisés par les agriculteurs camerounais.
\end{savaistu}

\courssub{La famille des halogènes : des non-métaux très réactifs}

\textbf{Présentation.} Les \cle{halogènes} occupent la 17\ieme{} colonne : fluor F, chlore Cl, brome Br, iode I. Ce sont des \textbf{non-métaux} très réactifs, qui aiment se combiner aux métaux pour former des sels (le mot « halogène » signifie d'ailleurs « qui engendre des sels »). Ainsi le chlore se combine au sodium pour donner le chlorure de sodium \ce{NaCl}, notre sel de table.

\textbf{Propriétés et usages.} Le chlore, corps simple \ce{Cl2}, est un gaz jaune-vert, âcre et \textbf{toxique}. Mais utilisé à petites doses, c'est un puissant \textbf{désinfectant} : il tue les microbes présents dans l'eau. C'est pourquoi on l'utilise pour traiter l'eau de boisson et pour fabriquer l'eau de Javel.

\begin{exemplebox}[Le chlore rend l'eau potable]
La \textbf{Camwater} (Cameroon Water Utilities), qui distribue l'eau potable à Douala, Yaoundé et dans d'autres villes du Cameroun, ajoute du chlore \ce{Cl2} à l'eau avant de l'envoyer dans les robinets. À faible dose, le chlore détruit les bactéries responsables du choléra et de la typhoïde, sans danger pour le consommateur. De même, l'\textbf{eau de Javel} que l'on utilise à la maison pour nettoyer et désinfecter contient des composés chlorés.
\end{exemplebox}

\begin{attention}[Danger : ne jamais mélanger eau de Javel et acide]
L'eau de Javel ne doit \textbf{jamais} être mélangée avec un acide (détartrant, vinaigre concentré, acide chlorhydrique). Ce mélange provoque un dégagement de \cle{dichlore} \ce{Cl2}, un \textbf{gaz toxique} qui attaque les voies respiratoires et peut être mortel dans une pièce mal aérée. À la maison, on utilise un seul produit à la fois, on rince abondamment, et on aère la pièce.
\end{attention}

\courssub{La famille des gaz nobles : les « paresseux » de la chimie}

\textbf{Présentation.} Les \cle{gaz nobles} occupent la 18\ieme{} et dernière colonne : hélium He, néon Ne, argon Ar, krypton Kr, xénon Xe. Ce sont des gaz \textbf{très peu réactifs} : on dit qu'ils sont \textbf{chimiquement inertes}, car ils ne se combinent pratiquement jamais avec d'autres éléments. Leur couche externe d'électrons est complète, ce qui les rend « satisfaits » et sans besoin de réagir.

\textbf{Applications.} Cette inertie est précieuse :
\begin{itemize}
\item le \textbf{néon} Ne, traversé par un courant électrique, émet une lumière rouge-orangée éclatante : c'est le gaz des \textbf{enseignes lumineuses} des boutiques et des bars de nos villes ;
\item l'\textbf{hélium} He, très léger et ininflammable, sert à gonfler les \textbf{ballons} (contrairement au dihydrogène, il ne risque pas d'exploser) ;
\item l'\textbf{argon} Ar remplit certaines lampes à incandescence et sert de gaz protecteur en soudure.
\end{itemize}

\begin{savaistu}[Le néon des enseignes et l'hélium des ballons]
Quand tu passes la nuit devant une enseigne lumineuse rougeoyante à Douala ou à Yaoundé, tu regardes du \textbf{néon} en train de briller ! Et le ballon qui s'envole à une fête est souvent rempli d'\textbf{hélium}. Ces deux gaz sont des gaz nobles : leur nom « noble » vient justement de leur « réserve » — comme les nobles d'autrefois, ils ne se « mélangent » pas avec les autres éléments.
\end{savaistu}

\courssub{Bilan et exercice d'application}

\begin{aretenir}[Les trois familles à connaître]
\begin{center}
\begin{tabularx}{\linewidth}{|l|X|X|X|}
\hline
\rowcolor{popblueL} \textbf{Famille} & \textbf{Colonne} & \textbf{Exemples} & \textbf{Caractère} \\
\hline
Alcalins & 1\iere{} (sauf H) & Li, Na, K & Métaux mous, réagissent violemment avec l'eau \\
\hline
Halogènes & 17\ieme{} & F, Cl, Br, I & Non-métaux très réactifs, désinfectants (chlore) \\
\hline
Gaz nobles & 18\ieme{} & He, Ne, Ar & Gaz inertes (très peu réactifs), lampes et enseignes \\
\hline
\end{tabularx}
\end{center}
\end{aretenir}

\begin{exoresolu}[Identifier une famille]
Le professeur place un petit morceau de potassium K dans un cristallisoir d'eau. La réaction est encore plus vive qu'avec le sodium : le potassium s'enflamme spontanément avec une flamme violette.\par
1. À quelle famille appartient le potassium ? Justifier.\par
2. Pourquoi ce résultat était-il prévisible sans faire l'expérience ?
\tcblower
\textbf{Solution.}\par
1. Le potassium K se trouve dans la \textbf{1\iere{} colonne} du tableau périodique : il appartient à la famille des \textbf{alcalins}, comme le sodium Na et le lithium Li.\par
2. Les éléments d'une même famille (même colonne) ont des \textbf{propriétés chimiques semblables}, car leurs atomes possèdent le même nombre d'électrons sur la couche externe. Sachant que le sodium réagit violemment avec l'eau, on pouvait donc prévoir que le potassium, son « frère » de famille, réagirait aussi violemment — et même davantage, car la réactivité des alcalins augmente en descendant la colonne.
\end{exoresolu}

\begin{attention}[Pièges fréquents au BEPC]
\begin{itemize}
\item Ne place \textbf{jamais} l'hydrogène H dans la famille des alcalins : il est dans la 1\iere{} colonne mais n'est \textbf{pas} un alcalin.
\item Ne confonds pas « famille » et « période » : une famille est une \textbf{colonne}, une période est une \textbf{ligne}.
\item Le chlore utilisé pour désinfecter est le corps simple \ce{Cl2} (dichlore) ; dans le sel \ce{NaCl}, il est combiné et inoffensif. Écris toujours les formules avec les bons indices.
\end{itemize}
\end{attention}

\begin{exercice}[À toi de jouer]
1. Recopie et complète : les éléments d'une même \ldots\ldots{} ont des propriétés chimiques semblables car leurs atomes ont le même nombre d'\ldots\ldots{} sur la couche externe.\par
2. Pourquoi conserve-t-on le sodium dans l'huile ?\par
3. Cite deux gaz nobles et donne une application de chacun.\par
4. Ta mère veut nettoyer les toilettes avec de l'eau de Javel puis un détartrant acide. Quel conseil lui donnes-tu et pourquoi ?
\end{exercice}

\begin{corrige}[Exercice « À toi de jouer »]
1. Les éléments d'une même \textbf{colonne (famille)} ont des propriétés chimiques semblables car leurs atomes ont le même nombre d'\textbf{électrons} sur la couche externe.\par
2. Le sodium réagit violemment avec l'eau et même avec l'humidité de l'air (dégagement de dihydrogène et de chaleur, risque d'inflammation). L'huile l'isole de l'air et de l'eau.\par
3. Exemples : le néon Ne, utilisé dans les enseignes lumineuses ; l'hélium He, utilisé pour gonfler les ballons (ou l'argon Ar, dans certaines lampes).\par
4. Je lui conseille de ne \textbf{jamais mélanger} les deux produits : le mélange eau de Javel + acide dégage du dichlore \ce{Cl2}, un gaz toxique dangereux pour les voies respiratoires. Il faut utiliser un seul produit, rincer abondamment et aérer la pièce.
\end{corrige}

\coursec{Éléments usuels et applications dans la vie courante}

Dans la section précédente, nous avons découvert quelques familles remarquables du tableau périodique : les métaux alcalins, les halogènes, les gaz nobles. Tu sais maintenant que le tableau de Mendeleïev range les éléments par numéro atomique croissant, en lignes (périodes) et en colonnes (familles). Mais à quoi sert ce grand tableau dans ta vie de tous les jours ? Réponse : à presque tout ! La marmite de maman, les tiges à béton du chantier voisin, le câble électrique d'ENEO, l'eau de Javel, et même ton propre sang contiennent des éléments du tableau. Dans cette dernière section, nous allons faire le tour des éléments usuels, apprendre à les repérer dans le tableau, et découvrir les règles de sécurité liées à certains produits de la maison.

\courssub{Les éléments métalliques usuels}

\begin{definition}[Métal usuel]
Un \cle{métal} est un élément chimique généralement solide à la température ambiante (sauf le mercure), brillant, bon conducteur de la chaleur et du courant électrique, et déformable sans se casser (on dit qu'il est \cle{malléable}). Dans le tableau périodique, les métaux occupent la partie gauche et le centre.
\end{definition}

Observe autour de toi : la plupart des objets fabriqués par l'homme contiennent des métaux. Voici les cinq métaux les plus importants de ton quotidien au Cameroun.

\textbf{Le fer (Fe, $Z = 26$).} C'est le métal de la \textbf{construction}. Les tiges à béton (fer à béton) que l'on voit sur tous les chantiers de Yaoundé et de Douala sont en fer. Le fer est dur et résistant, mais il rouille au contact de l'air humide : c'est pourquoi on le protège par de la peinture. Le fer est aussi un élément essentiel de ton sang (nous y reviendrons).

\textbf{L'aluminium (Al, $Z = 13$).} Léger et inoxydable, l'aluminium sert à fabriquer les \textbf{marmites}, les casseroles, les canettes de boisson et les tôles. C'est un excellent conducteur de la chaleur : la marmite en aluminium chauffe vite et uniformément.

\textbf{Le cuivre (Cu, $Z = 29$).} C'est le métal de l'\textbf{électricité}. Les câbles électriques des maisons et les lignes d'ENEO contiennent du cuivre, car c'est un excellent conducteur du courant, juste derrière l'argent (trop cher). Le vol des câbles de cuivre est d'ailleurs un vrai problème : il prive des quartiers entiers d'électricité.

\textbf{Le zinc (Zn, $Z = 30$).} Les \textbf{tôles galvanisées} qui couvrent beaucoup de maisons sont des tôles d'acier recouvertes d'une fine couche de zinc. Le zinc protège l'acier de la rouille : c'est la galvanisation.

\textbf{L'or (Au, $Z = 79$).} Métal précieux qui ne s'oxyde jamais, l'or sert à fabriquer les \textbf{bijoux} et les alliances. Le Cameroun possède des gisements d'or, notamment dans la région de l'Est.

\begin{exemplebox}[L'aluminium dans le tableau]
L'aluminium a pour symbole \ce{Al} et pour numéro atomique $Z = 13$. Dans le tableau périodique, on le trouve à la \textbf{3\textsuperscript{e} période} (3\textsuperscript{e} ligne) et dans la \textbf{13\textsuperscript{e} colonne}. C'est un \cle{métal} : il est brillant, conduit la chaleur (marmites) et le courant électrique. Sa légèreté (il est environ trois fois moins dense que le fer) explique son usage dans les canettes et l'aviation.
\end{exemplebox}

\courssub{Les éléments non métalliques usuels}

Les non-métaux occupent la partie droite du tableau périodique. Ils sont moins visibles que les métaux, mais tout aussi indispensables.

\textbf{L'oxygène (O, $Z = 8$).} Le dioxygène \ce{O2} de l'air est indispensable à la \textbf{respiration} : sans lui, pas de vie possible. À l'hôpital, les malades en détresse respiratoire reçoivent de l'oxygène médical stocké dans des bouteilles.

\textbf{Le carbone (C, $Z = 6$).} Le carbone est l'élément de base de tous les \textbf{êtres vivants} : nos muscles, nos os, nos cheveux contiennent du carbone. On le trouve aussi dans le charbon de bois utilisé pour la cuisine, dans le graphite des crayons et dans le diamant.

\textbf{L'azote (N, $Z = 7$).} Le diazote \ce{N2} forme environ 78 \% de l'air. L'azote entre dans la composition des \textbf{engrais} qui nourrissent les plantes : les agriculteurs de l'Ouest l'utilisent pour améliorer les récoltes de maïs et de pommes de terre.

\textbf{Le soufre (S, $Z = 16$).} Solide jaune à l'odeur caractéristique, le soufre sert à fabriquer l'acide sulfurique, des médicaments et des produits de traitement des plantes.

\textbf{Le chlore (Cl, $Z = 17$).} Le chlore est utilisé pour la \textbf{désinfection} : l'eau de Javel et l'eau de forage traitée contiennent des composés chlorés qui tuent les microbes. La Camwater traite l'eau potable avec du chlore.

\begin{popfigure}
\begin{center}
\begin{tikzpicture}[mindmap, grow cyclic,
  every node/.style={concept, font=\small},
  root concept/.append style={concept color=popblue, font=\large\bfseries, text=white},
  level 1/.append style={level distance=3.6cm, sibling angle=72},
  level 1 concept/.append style={concept color=popblueL, text=popdark},
  level 2/.append style={level distance=2.6cm, sibling angle=45},
  level 2 concept/.append style={concept color=poporangeL, text=popdark, font=\footnotesize}]
  \node[concept] {Éléments usuels}
    child { node[concept] {Fe (fer)} child { node[concept] {tiges à béton} } }
    child { node[concept] {Al (aluminium)} child { node[concept] {marmites, canettes} } }
    child { node[concept] {Cu (cuivre)} child { node[concept] {câbles électriques} } }
    child { node[concept] {Cl (chlore)} child { node[concept] {eau de Javel} } }
    child { node[concept] {Ca (calcium)} child { node[concept] {os et dents} } };
\end{tikzpicture}
\end{center}
\legende{Carte mentale : chaque élément usuel est relié à une application concrète de la vie courante au Cameroun.}
\end{popfigure}

\courssub{Les éléments essentiels au vivant}

Ton corps est une véritable « usine chimique » construite à partir d'une vingtaine d'éléments du tableau périodique. Quatre éléments forment à eux seuls environ 96 \% de la masse du corps humain.

\begin{center}
\begin{tabularx}{\linewidth}{|l|c|X|}
\hline
\rowcolor{popblueL} \textbf{Élément} & \textbf{Symbole ($Z$)} & \textbf{Rôle dans le corps} \\
\hline
Oxygène & O (8) & Respiration, composition de l'eau du corps \\
\hline
Carbone & C (6) & Base de toutes les molécules du vivant \\
\hline
Hydrogène & H (1) & Composition de l'eau et des molécules du vivant \\
\hline
Azote & N (7) & Composition des protéines (muscles) \\
\hline
Calcium & Ca (20) & Solidité des os et des dents \\
\hline
Fer & Fe (26) & Transport de l'oxygène par le sang (hémoglobine) \\
\hline
Potassium & K (19) & Fonctionnement des cellules et des nerfs \\
\hline
Sodium & Na (11) & Fonctionnement des cellules, équilibre de l'eau \\
\hline
\end{tabularx}
\end{center}

\begin{aretenir}[À retenir]
Les quatre éléments \ce{C}, \ce{H}, \ce{O}, \ce{N} constituent l'essentiel de la matière vivante. Le \cle{calcium} \ce{Ca} solidifie les os, le \cle{fer} \ce{Fe} permet au sang de transporter le dioxygène, et les ions \cle{potassium} \ce{K+} et \cle{sodium} \ce{Na+} assurent le fonctionnement électrique des cellules (nerfs, muscles, cœur).
\end{aretenir}

\begin{savaistu}[Le fer dans ton assiette]
L'hémoglobine, le pigment rouge du sang, contient un atome de \textbf{fer} en son centre : c'est lui qui accroche le dioxygène dans les poumons pour le distribuer à tout le corps. Quand on manque de fer, on souffre d'\textbf{anémie} : fatigue, pâleur, essoufflement. C'est pourquoi les médecins conseillent de manger des feuilles de \textbf{folong}, du \textbf{ndolé} et du \textbf{zom}, des légumes camerounais très riches en fer, ainsi que de la viande et du poisson. Grandir avec assez de fer, c'est aussi mieux réussir à l'école !
\end{savaistu}

\courssub{Lecture pratique du tableau périodique}

Face à un exercice du BEPC, on te demandera souvent de repérer un élément dans le tableau : son numéro atomique, sa période, sa famille. Voici la méthode infaillible.

\begin{methode}[Exploiter le tableau périodique]
Pour répondre à une question sur un élément :
\begin{enumerate}
\item \textbf{Repère le symbole} de l'élément dans le tableau (une ou deux lettres, la première en majuscule : \ce{Fe}, \ce{Al}, \ce{Cl}...).
\item \textbf{Lis le numéro atomique $Z$} : c'est le nombre entier écrit dans la case, en général au-dessus du symbole. Il indique le nombre de protons (et d'électrons pour l'atome neutre).
\item \textbf{Compte la ligne} (période) : elle donne le nombre de couches électroniques occupées.
\item \textbf{Compte la colonne} (famille) : les éléments d'une même colonne ont des propriétés chimiques voisines.
\item \textbf{Conclus} : précise si l'élément est un métal (gauche et centre) ou un non-métal (droite).
\end{enumerate}
\end{methode}

\begin{exoresolu}[Repérer le chlore dans le tableau]
Énoncé. Le chlore \ce{Cl} est utilisé pour désinfecter l'eau de Javel. Donne son numéro atomique, sa période et sa famille. Précise s'il s'agit d'un métal ou d'un non-métal.
\tcblower
Solution. On applique la méthode pas à pas.
\begin{enumerate}
\item On repère le symbole \ce{Cl} dans le tableau périodique.
\item Le nombre écrit dans la case est $Z = 17$ : l'atome de chlore possède 17 protons et 17 électrons.
\item Le chlore se trouve sur la \textbf{3\textsuperscript{e} ligne} : il appartient à la 3\textsuperscript{e} période (3 couches électroniques).
\item Il se trouve dans la \textbf{17\textsuperscript{e} colonne} : c'est la famille des \textbf{halogènes}, comme le fluor \ce{F}.
\item Le chlore est à droite du tableau : c'est un \textbf{non-métal}.
\end{enumerate}
Conclusion : le chlore \ce{Cl} ($Z = 17$) est un non-métal de la 3\textsuperscript{e} période, de la famille des halogènes.
\end{exoresolu}

\begin{attention}[Un élément n'est pas une molécule !]
Erreur très fréquente au BEPC : confondre \cle{élément} et \cle{composé}. Le fer \ce{Fe} est un \textbf{élément} : il est formé d'un seul type d'atome. L'eau \ce{H2O} est un \textbf{composé} : sa molécule est formée de \textbf{deux éléments différents} (hydrogène et oxygène) chimiquement liés. De même, le sel de cuisine \ce{NaCl} est un composé de deux éléments : le sodium et le chlore. Un élément ne peut pas être « décomposé » en substances plus simples ; un composé, si.
\end{attention}

\courssub{Sécurité domestique : lire les étiquettes}

Certains produits de la maison contiennent des éléments ou des composés dangereux : l'eau de Javel (composés chlorés), les déboucheurs de canalisations (soude ou acides concentrés), les acides de batterie. Pour protéger les usagers, les fabricants apposent des \cle{pictogrammes de danger} sur les étiquettes : des losanges rouges à fond blanc contenant un symbole noir.

\begin{popfigure}
\begin{center}
\begin{tikzpicture}[scale=1]
% Pictogramme corrosif
\begin{scope}[shift={(0,0)}]
  \draw[fill=white, draw=popink, line width=1.2pt, rotate=45] (-1.15,-1.15) rectangle (1.15,1.15);
  \draw[line width=1pt] (-0.55,0.55) -- (-0.1,0.1) -- (-0.1,-0.35);
  \draw[line width=1pt] (0.55,0.55) -- (0.1,0.1) -- (0.1,-0.35);
  \draw[line width=1.4pt] (-0.35,-0.5) -- (0.35,-0.5);
  \draw[line width=1pt] (-0.1,-0.35) -- (-0.25,-0.5);
  \draw[line width=1pt] (0.1,-0.35) -- (0.25,-0.5);
  \node[below, font=\small\bfseries] at (0,-1.5) {Corrosif};
\end{scope}
% Pictogramme toxique
\begin{scope}[shift={(4.2,0)}]
  \draw[fill=white, draw=popink, line width=1.2pt, rotate=45] (-1.15,-1.15) rectangle (1.15,1.15);
  \draw[fill=popdark] (0,0.25) circle (0.28);
  \draw[fill=white] (-0.11,0.33) circle (0.07);
  \draw[fill=white] (0.11,0.33) circle (0.07);
  \draw[fill=white] (0,0.13) circle (0.05);
  \draw[line width=1.4pt] (-0.35,-0.35) -- (0.35,-0.15);
  \draw[line width=1.4pt] (-0.35,-0.15) -- (0.35,-0.35);
  \fill[popdark] (-0.35,-0.35) circle (0.06);
  \fill[popdark] (0.35,-0.15) circle (0.06);
  \fill[popdark] (-0.35,-0.15) circle (0.06);
  \fill[popdark] (0.35,-0.35) circle (0.06);
  \node[below, font=\small\bfseries] at (0,-1.5) {Toxique};
\end{scope}
\end{tikzpicture}
\end{center}
\legende{Deux pictogrammes de danger présents sur les produits ménagers : corrosif (attaque la peau et les métaux, ex. déboucheur) et toxique (dangereux en cas d'ingestion ou d'inhalation).}
\end{popfigure}

\begin{aretenir}[Consignes de sécurité à la maison]
\begin{itemize}
\item Toujours \textbf{lire l'étiquette} avant d'utiliser un produit ménager.
\item Ne \textbf{jamais mélanger} l'eau de Javel avec un acide (déboucheur, vinaigre concentré) : il se dégage un gaz toxique, le dichlore \ce{Cl2}.
\item Porter des \textbf{gants} pour manipuler les produits corrosifs et garder ces produits \textbf{hors de portée des enfants}.
\item En cas de contact avec la peau ou les yeux : rincer abondamment à l'eau claire et consulter un médecin.
\end{itemize}
\end{aretenir}

\begin{exercice}[À toi de jouer]
\begin{enumerate}
\item Le zinc \ce{Zn} a pour numéro atomique $Z = 30$. À quoi sert-il dans la construction des maisons ?
\item Cite trois éléments essentiels au corps humain et précise le rôle de chacun.
\item Pourquoi ne faut-il jamais mélanger l'eau de Javel avec un déboucheur acide ?
\end{enumerate}
\end{exercice}

\begin{corrige}[Exercice : éléments usuels et sécurité]
\begin{enumerate}
\item Le zinc sert à recouvrir les tôles d'acier : on obtient des \textbf{tôles galvanisées} protégées de la rouille, utilisées pour les toitures.
\item Le \textbf{calcium} \ce{Ca} solidifie les os et les dents ; le \textbf{fer} \ce{Fe} permet à l'hémoglobine du sang de transporter le dioxygène ; le \textbf{sodium} \ce{Na} (ou le potassium \ce{K}) assure le fonctionnement des cellules, des nerfs et des muscles.
\item Ce mélange dégage du \textbf{dichlore} \ce{Cl2}, un gaz toxique qui irrite gravement les poumons et les yeux. Il faut donc utiliser ces produits séparément et lire les étiquettes.
\end{enumerate}
\end{corrige}

\courssub{Bilan de la leçon}

\begin{aretenir}[Le tableau périodique : un outil de rangement]
Le tableau de Mendeleïev n'est pas qu'un poster à apprendre par cœur : c'est un \textbf{outil de rangement et de prévision}. En rangeant les éléments par numéro atomique croissant, il regroupe dans une même colonne des éléments aux propriétés voisines (même famille) et permet de prévoir le comportement d'un élément que l'on n'a jamais étudié. Connaître la place d'un élément, c'est déjà connaître une partie de ses propriétés : métal ou non-métal, réactif ou inerte, utile ou dangereux. De la marmite en aluminium au fer de ton sang, tout ton quotidien est écrit dans ce tableau !
\end{aretenir}

\newpage

\coursec{Activité d'intégration}

\courssub{Objectifs de l'activité}

À la fin de cette activité, tu seras capable de :
\begin{itemize}
\item identifier l'\cle{élément chimique} principal contenu dans un objet ou un produit de la vie courante ;
\item écrire correctement le \cle{symbole} d'un élément et relever son \cle{numéro atomique} $Z$ dans le tableau périodique ;
\item situer un élément dans le tableau de Mendeleïev : numéro de la \cle{période} (ligne) et nom de la \cle{famille} (colonne) ;
\item distinguer un \cle{métal} d'un \cle{non-métal} à partir de sa position dans le tableau ;
\item appliquer une consigne de sécurité en la justifiant par les propriétés chimiques d'une famille d'éléments.
\end{itemize}

\courssub{Situation}

Pendant les vacances, le jeune Tabi, élève de 3ème au collège de Nkongsamba, accompagne sa mère au grand marché de la ville. Dans les différents rayons, il observe les objets et produits que les commerçants vendent : des marmites en aluminium brillantes, des tiges à béton pour la construction, des rouleaux de câble électrique en cuivre, des flacons d'eau de Javel, des sacs de sel de cuisine, des sacs de charbon de bois, des bonbonnes de gaz butane, des sacs d'engrais NPK pour les champs de bananiers, des boîtes de comprimés de fer vendues à la pharmacie du marché, et des tôles de zinc pour la toiture des maisons.

Tabi se souvient de la leçon sur le tableau de classification de Mendeleïev. Il se pose alors les questions suivantes : \emph{quels éléments chimiques se cachent derrière tous ces produits ? Où se situent-ils dans le tableau périodique ? Sont-ce des métaux ou des non-métaux ?}

De retour en classe, l'enseignant de PCT transforme cette promenade au marché en enquête scientifique. Il distribue à chaque groupe d'élèves un \textbf{tableau périodique simplifié} (les 20 premiers éléments, plus quelques éléments usuels comme Fe, Cu, Zn, Br, I) et la liste des 10 produits observés par Tabi, avec les indications suivantes :

\begin{center}
\begin{tabularx}{\linewidth}{|c|X|l|}
\hline
\rowcolor{popblueL} \textbf{N°} & \textbf{Objet ou produit} & \textbf{Indice chimique} \\
\hline
1 & Marmite en aluminium & Corps simple pur \\
\hline
2 & Tige à béton (fer à béton) & Corps simple métallique \\
\hline
3 & Câble électrique (âme en cuivre) & Corps simple métallique \\
\hline
4 & Eau de Javel & Contient l'ion hypochlorite \ce{ClO-} \\
\hline
5 & Sel de cuisine & Chlorure de sodium \ce{NaCl} \\
\hline
6 & Charbon de bois & Presque uniquement du carbone \\
\hline
7 & Bonbonne de gaz butane & \ce{C4H10} : carbone et hydrogène \\
\hline
8 & Engrais NPK & Azote N, phosphore P, potassium K \\
\hline
9 & Comprimé de fer & Contient l'élément fer \\
\hline
10 & Tôle de zinc (toiture) & Corps simple métallique \\
\hline
\end{tabularx}
\end{center}

Sur le flacon d'eau de Javel, Tabi lit la mention : \og Ne pas mélanger avec des produits acides (détartrants, vinaigre) : dégagement d'un gaz toxique \fg. La classe devra expliquer cette consigne à la fin de l'enquête.

\courssub{Consignes}

\textbf{Travail de groupe (30 min), puis présentation orale d'un cas par groupe.}

\begin{enumerate}
\item Pour chacun des 10 produits, identifie l'\cle{élément chimique} principal (pour le gaz butane, choisis le carbone ; pour l'engrais NPK, choisis le potassium).
\item Écris le \cle{symbole} de chaque élément en respectant la règle d'écriture (une majuscule, éventuellement suivie d'une minuscule).
\item À l'aide du tableau périodique simplifié, relève le \cle{numéro atomique} $Z$ de chaque élément.
\item Indique le numéro de la \cle{période} (ligne) et le nom de la \cle{famille} (colonne) de chaque élément.
\item Précise si l'élément est un \cle{métal} ou un \cle{non-métal}, en justifiant par sa position dans le tableau (à gauche et au centre : métaux ; à droite : non-métaux).
\item Consigne tous les résultats dans un tableau à double entrée à 6 colonnes : produit, élément, symbole, $Z$, période/famille, métal ou non-métal.
\item \textbf{Question de synthèse :} le chlore appartient à la famille des halogènes (colonne 17). Explique pourquoi l'eau de Javel ne doit jamais être mélangée avec un acide, en t'appuyant sur la réactivité des halogènes et sur le gaz \ce{Cl2} qui peut se dégager.
\end{enumerate}

\courssub{Ressources à mobiliser}

\begin{aretenir}[Les outils de la leçon]
\begin{itemize}
\item Le \cle{symbole} d'un élément commence toujours par une \textbf{majuscule}, suivie éventuellement d'une minuscule : \ce{Fe}, \ce{Na}, \ce{Cl}.
\item Le \cle{numéro atomique} $Z$ est le nombre de protons du noyau ; c'est aussi le numéro de la case de l'élément dans le tableau.
\item Les \textbf{lignes} du tableau sont les \cle{périodes} ; les \textbf{colonnes} sont les \cle{familles} : alcalins (colonne 1), alcalino-terreux (colonne 2), halogènes (colonne 17), gaz nobles (colonne 18).
\item Les éléments situés à gauche et au centre du tableau sont des \cle{métaux} ; ceux situés à droite sont des \cle{non-métaux}.
\item Les éléments d'une même famille ont des \textbf{propriétés chimiques voisines}.
\end{itemize}
\end{aretenir}

\begin{popfigure}
\begin{center}
\begin{tikzpicture}[scale=0.62, every node/.style={font=\small}]
% Période 1
\draw[fill=popgreenL] (0,0) rectangle (1,1); \node at (0.5,0.6) {\textbf{H}}; \node at (0.5,0.25) {\tiny 1};
\draw[fill=popgreenL] (17,0) rectangle (18,1); \node at (17.5,0.6) {\textbf{He}}; \node at (17.5,0.25) {\tiny 2};
% Période 2
\draw[fill=poporangeL] (0,-1) rectangle (1,0); \node at (0.5,-0.4) {\textbf{Li}}; \node at (0.5,-0.75) {\tiny 3};
\draw[fill=poporangeL] (1,-1) rectangle (2,0); \node at (1.5,-0.4) {\textbf{Be}}; \node at (1.5,-0.75) {\tiny 4};
\draw[fill=popgreenL] (12,-1) rectangle (13,0); \node at (12.5,-0.4) {\textbf{B}}; \node at (12.5,-0.75) {\tiny 5};
\draw[fill=popgreenL] (13,-1) rectangle (14,0); \node at (13.5,-0.4) {\textbf{C}}; \node at (13.5,-0.75) {\tiny 6};
\draw[fill=popgreenL] (14,-1) rectangle (15,0); \node at (14.5,-0.4) {\textbf{N}}; \node at (14.5,-0.75) {\tiny 7};
\draw[fill=popgreenL] (15,-1) rectangle (16,0); \node at (15.5,-0.4) {\textbf{O}}; \node at (15.5,-0.75) {\tiny 8};
\draw[fill=popgreenL] (16,-1) rectangle (17,0); \node at (16.5,-0.4) {\textbf{F}}; \node at (16.5,-0.75) {\tiny 9};
\draw[fill=popgreenL] (17,-1) rectangle (18,0); \node at (17.5,-0.4) {\textbf{Ne}}; \node at (17.5,-0.75) {\tiny 10};
% Période 3
\draw[fill=poporangeL] (0,-2) rectangle (1,-1); \node at (0.5,-1.4) {\textbf{Na}}; \node at (0.5,-1.75) {\tiny 11};
\draw[fill=poporangeL] (1,-2) rectangle (2,-1); \node at (1.5,-1.4) {\textbf{Mg}}; \node at (1.5,-1.75) {\tiny 12};
\draw[fill=poporangeL] (12,-2) rectangle (13,-1); \node at (12.5,-1.4) {\textbf{Al}}; \node at (12.5,-1.75) {\tiny 13};
\draw[fill=popgreenL] (13,-2) rectangle (14,-1); \node at (13.5,-1.4) {\textbf{Si}}; \node at (13.5,-1.75) {\tiny 14};
\draw[fill=popgreenL] (14,-2) rectangle (15,-1); \node at (14.5,-1.4) {\textbf{P}}; \node at (14.5,-1.75) {\tiny 15};
\draw[fill=popgreenL] (15,-2) rectangle (16,-1); \node at (15.5,-1.4) {\textbf{S}}; \node at (15.5,-1.75) {\tiny 16};
\draw[fill=popgreenL] (16,-2) rectangle (17,-1); \node at (16.5,-1.4) {\textbf{Cl}}; \node at (16.5,-1.75) {\tiny 17};
\draw[fill=popgreenL] (17,-2) rectangle (18,-1); \node at (17.5,-1.4) {\textbf{Ar}}; \node at (17.5,-1.75) {\tiny 18};
% Période 4
\draw[fill=poporangeL] (0,-3) rectangle (1,-2); \node at (0.5,-2.4) {\textbf{K}}; \node at (0.5,-2.75) {\tiny 19};
\draw[fill=poporangeL] (1,-3) rectangle (2,-2); \node at (1.5,-2.4) {\textbf{Ca}}; \node at (1.5,-2.75) {\tiny 20};
\draw[fill=poporangeL] (7,-3) rectangle (8,-2); \node at (7.5,-2.4) {\textbf{Fe}}; \node at (7.5,-2.75) {\tiny 26};
\draw[fill=poporangeL] (10,-3) rectangle (11,-2); \node at (10.5,-2.4) {\textbf{Cu}}; \node at (10.5,-2.75) {\tiny 29};
\draw[fill=poporangeL] (11,-3) rectangle (12,-2); \node at (11.5,-2.4) {\textbf{Zn}}; \node at (11.5,-2.75) {\tiny 30};
\draw[fill=popgreenL] (16,-3) rectangle (17,-2); \node at (16.5,-2.4) {\textbf{Br}}; \node at (16.5,-2.75) {\tiny 35};
% Période 5 (iode)
\draw[fill=popgreenL] (16,-4) rectangle (17,-3); \node at (16.5,-3.4) {\textbf{I}}; \node at (16.5,-3.75) {\tiny 53};
% Annotations familles
\node[font=\footnotesize\bfseries, popdark] at (0.5,1.4) {Alcalins};
\node[font=\footnotesize\bfseries, popdark] at (1.5,1.4) {Alcalino-terreux};
\node[font=\footnotesize\bfseries, popdark] at (16.5,1.4) {Halogènes};
\node[font=\footnotesize\bfseries, popdark] at (17.5,1.4) {Gaz nobles};
% Annotations périodes
\node[font=\footnotesize\bfseries, popdark, rotate=90] at (-0.7,0.5) {P. 1};
\node[font=\footnotesize\bfseries, popdark, rotate=90] at (-0.7,-0.5) {P. 2};
\node[font=\footnotesize\bfseries, popdark, rotate=90] at (-0.7,-1.5) {P. 3};
\node[font=\footnotesize\bfseries, popdark, rotate=90] at (-0.7,-2.5) {P. 4};
% Légende couleurs
\draw[fill=poporangeL] (4,-5.2) rectangle (4.6,-4.6); \node[right, font=\footnotesize] at (4.7,-4.9) {Métaux};
\draw[fill=popgreenL] (8,-5.2) rectangle (8.6,-4.6); \node[right, font=\footnotesize] at (8.7,-4.9) {Non-métaux};
\end{tikzpicture}
\end{center}
\legende{Extrait simplifié du tableau périodique utilisé pour l'enquête. Dans chaque case : le symbole de l'élément et, en dessous, son numéro atomique $Z$. Les lignes horizontales sont les périodes (P. 1 à P. 4), les colonnes verticales sont les familles. Les métaux sont en orange, les non-métaux en vert.}
\end{popfigure}

\courssub{Démarche et corrigé commenté}

\begin{exoresolu}[Enquête au marché : corrigé complet]
Énoncé : identifier et classer dans le tableau périodique les éléments présents dans les 10 produits du marché de Nkongsamba, puis justifier la consigne de sécurité de l'eau de Javel.

\tcblower

\textbf{Étape 1 : identification de l'élément principal.} On utilise les indices chimiques et les connaissances de la vie courante : la marmite est en aluminium ; le fer à béton et le comprimé contiennent du fer ; l'âme du câble est en cuivre ; l'eau de Javel contient du chlore ; le sel contient du sodium et du chlore (on retient le sodium, élément caractéristique) ; le charbon est du carbone ; le butane contient du carbone ; l'engrais contient du potassium ; la tôle est en zinc.

\textbf{Étape 2 : lecture du tableau périodique.} Pour chaque élément, on relève le symbole, le numéro atomique $Z$, la période et la famille. On obtient le tableau à double entrée suivant :

\begin{center}
\begin{tabularx}{\linewidth}{|X|l|c|c|X|}
\hline
\rowcolor{popblueL} \textbf{Produit} & \textbf{Élément (symbole)} & $Z$ & \textbf{Période} & \textbf{Famille / nature} \\
\hline
Marmite & Aluminium \ce{Al} & 13 & 3 & Métal (colonne 13) \\
\hline
Tige à béton & Fer \ce{Fe} & 26 & 4 & Métal de transition \\
\hline
Câble électrique & Cuivre \ce{Cu} & 29 & 4 & Métal de transition \\
\hline
Eau de Javel & Chlore \ce{Cl} & 17 & 3 & Halogène, non-métal \\
\hline
Sel de cuisine & Sodium \ce{Na} & 11 & 3 & Alcalin, métal \\
\hline
Charbon de bois & Carbone \ce{C} & 6 & 2 & Non-métal (colonne 14) \\
\hline
Gaz butane & Carbone \ce{C} & 6 & 2 & Non-métal (colonne 14) \\
\hline
Engrais NPK & Potassium \ce{K} & 19 & 4 & Alcalin, métal \\
\hline
Comprimé & Fer \ce{Fe} & 26 & 4 & Métal de transition \\
\hline
Tôle & Zinc \ce{Zn} & 30 & 4 & Métal de transition \\
\hline
\end{tabularx}
\end{center}

\textbf{Étape 3 : vérifications.} Le symbole s'écrit avec une majuscule puis une minuscule : on écrit \ce{Cl} et non CL, \ce{Fe} et non FE. Le numéro atomique croît quand on parcourt le tableau de gauche à droite et de haut en bas : \ce{C} (6) $<$ \ce{Na} (11) $<$ \ce{Al} (13) $<$ \ce{Cl} (17) $<$ \ce{K} (19) $<$ \ce{Fe} (26) $<$ \ce{Cu} (29) $<$ \ce{Zn} (30). Les métaux (\ce{Al}, \ce{Fe}, \ce{Cu}, \ce{Na}, \ce{K}, \ce{Zn}) sont bien situés à gauche et au centre ; les non-métaux (\ce{C}, \ce{Cl}) à droite.

\textbf{Étape 4 : question de synthèse (sécurité).} Le chlore \ce{Cl} appartient à la famille des \cle{halogènes} (colonne 17), des non-métaux très réactifs. L'eau de Javel contient des ions hypochlorite \ce{ClO-}. Si on la mélange avec un acide (détartrant, vinaigre), une réaction chimique dégage du \textbf{dichlore} \ce{Cl2}, un gaz toxique qui attaque les voies respiratoires. La réactivité élevée des halogènes, commune à toute la famille (F, Cl, Br, I), explique ce danger. Consigne : ne jamais mélanger l'eau de Javel avec un acide, et l'utiliser dans un endroit aéré.
\end{exoresolu}

\begin{attention}[Erreurs fréquentes à éviter]
\begin{itemize}
\item Ne confonds pas le \textbf{symbole} (\ce{Fe}) avec le nom de l'élément (fer) ni avec la formule d'un composé (\ce{NaCl} est le sel, pas un élément).
\item Le numéro atomique $Z$ n'est pas la masse : c'est le numéro de la case dans le tableau.
\item \og Famille \fg{} désigne la \textbf{colonne}, \og période \fg{} la \textbf{ligne} : ne les inverse pas dans ta rédaction au BEPC.
\item Le sodium \ce{Na} et le potassium \ce{K} sont des métaux même si on ne les voit jamais sous forme de tôle : ils sont trop réactifs pour exister à l'état pur dans la nature.
\end{itemize}
\end{attention}

\courssub{Bilan de l'activité}

Cette enquête au marché a montré que le tableau de Mendeleïev n'est pas un simple poster de classe : c'est une véritable \textbf{carte d'identité de la matière} qui nous entoure. En partant d'objets familiers du quotidien camerounais (marmite, tôle, bonbonne de gaz, engrais), tu as appris à remonter de l'objet à l'élément, puis de l'élément à sa case dans le tableau : symbole, numéro atomique $Z$, période et famille.

L'activité a aussi mis en évidence deux idées essentielles. D'abord, la \textbf{position d'un élément dans le tableau renseigne sur sa nature} : les métaux (fer, cuivre, zinc, aluminium) servent à fabriquer les objets solides et conducteurs, tandis que les non-métaux (carbone, chlore) entrent dans les combustibles et les produits chimiques. Ensuite, les éléments d'une \textbf{même famille ont des propriétés voisines}, ce qui permet de prévoir des dangers : la réactivité des halogènes explique la consigne de sécurité inscrite sur les flacons d'eau de Javel. Classer, c'est donc comprendre et se protéger.

\newpage

\coursec{Méthodes et exercices résolus}

\coursec{Classification des éléments}

\courssub{Rappels : éléments chimiques et symboles}

\begin{definition}[Élément chimique]
Un \cle{élément chimique} est une espèce chimique caractérisée par son \cle{numéro atomique} $Z$ (nombre de protons dans le noyau). Il ne peut pas être décomposé en substances plus simples par des réactions chimiques ordinaires. Exemples : le fer \ce{Fe}, l'oxygène \ce{O}, le carbone \ce{C}.
\end{definition}

\begin{definition}[Symbole chimique]
Le \cle{symbole} d'un élément s'écrit avec \textbf{une lettre majuscule} éventuellement suivie \textbf{d'une lettre minuscule}. Il est souvent issu du nom latin ou français.
\begin{center}
\begin{tabularx}{\linewidth}{|l|X|X|}\hline
\rowcolor{popblueL} Élément & Symbole & Origine \\ \hline
Hydrogène & \ce{H} & Hydrogène \\ \hline
Carbone & \ce{C} & Carbone \\ \hline
Sodium & \ce{Na} & \emph{Natrium} (latin) \\ \hline
Fer & \ce{Fe} & \emph{Ferrum} (latin) \\ \hline
Chlore & \ce{Cl} & Chlore \\ \hline
Potassium & \ce{K} & \emph{Kalium} (latin) \\ \hline
\end{tabularx}
\end{center}
\end{definition}

\begin{methode}[Retrouver un élément dans le tableau de Mendeleïev]
Pour identifier un élément chimique dont on connaît le symbole (ou l'inverse) :
\begin{enumerate}
\item Bien écrire le symbole : une \cle{lettre majuscule}, éventuellement suivie d'une \cle{minuscule} (ex. \ce{Fe}, \ce{Cl}, \ce{Na}).
\item Repérer dans le tableau la case portant ce symbole.
\item Lire le \cle{nom} de l'élément et son \cle{numéro atomique} $Z$ (le nombre entier écrit dans la case).
\item Vérifier la cohérence : le symbole dérive souvent du nom latin (ex. \ce{Na} pour \emph{natrium} = sodium, \ce{Fe} pour \emph{ferrum} = fer, \ce{K} pour \emph{kalium} = potassium).
\end{enumerate}
\textbf{Réflexe :} apprendre par cœur les symboles des éléments usuels : \ce{H}, \ce{C}, \ce{N}, \ce{O}, \ce{Na}, \ce{Mg}, \ce{Al}, \ce{S}, \ce{Cl}, \ce{K}, \ce{Ca}, \ce{Fe}, \ce{Cu}, \ce{Zn}, \ce{Ag}, \ce{Au}, \ce{Pb}.
\end{methode}

\begin{exoresolu}[Symboles et noms des éléments]
\textbf{Énoncé.} Compléter le tableau suivant :
\begin{center}
\begin{tabularx}{\linewidth}{|l|X|X|}\hline
\rowcolor{popblueL} Symbole & Nom de l'élément & Numéro atomique $Z$ \\ \hline
\ce{Na} & \ldots & \ldots \\ \hline
\ldots & Chlore & \ldots \\ \hline
\ce{Fe} & \ldots & \ldots \\ \hline
\ldots & \ldots & 8 \\ \hline
\end{tabularx}
\end{center}
\tcblower
\textbf{Solution.}
\begin{enumerate}
\item \ce{Na} : la case du tableau portant le symbole \ce{Na} correspond au \cle{sodium}, de numéro atomique $Z = 11$.
\item Chlore : son symbole est \ce{Cl} (attention : C majuscule, l minuscule), $Z = 17$.
\item \ce{Fe} : c'est le \cle{fer} (du latin \emph{ferrum}), $Z = 26$.
\item L'élément de numéro atomique $Z = 8$ est l'\cle{oxygène}, de symbole \ce{O}.
\end{enumerate}
\begin{center}
\begin{tabularx}{\linewidth}{|l|X|X|}\hline
\rowcolor{popblueL} Symbole & Nom de l'élément & Numéro atomique $Z$ \\ \hline
\ce{Na} & Sodium & 11 \\ \hline
\ce{Cl} & Chlore & 17 \\ \hline
\ce{Fe} & Fer & 26 \\ \hline
\ce{O} & Oxygène & 8 \\ \hline
\end{tabularx}
\end{center}
\textbf{Conclusion :} chaque élément est repéré de façon unique par son symbole et son numéro atomique $Z$.
\end{exoresolu}

\courssub{Historique et construction du tableau de Mendeleïev}

\begin{savaistu}[Dmitri Mendeleïev, l'architecte du tableau]
En 1869, le chimiste russe \textbf{Dmitri Mendeleïev} classe les 63 éléments connus par \textbf{masse atomique croissante}. Il observe une \cle{périodicité} des propriétés : les éléments aux propriétés semblables reviennent à intervalles réguliers. Il a l'audace de \textbf{laisser des cases vides} pour des éléments encore inconnus (ex. \emph{eka-aluminium}, \emph{eka-silicium}) et \textbf{prédit leurs propriétés} avec une précision stupéfiante. La découverte du gallium (1875), du scandium (1879) et du germanium (1886) confirme sa classification. En 1913, Henry Moseley montre que l'ordre fondamental est le \textbf{numéro atomique $Z$} (nombre de protons), et non la masse atomique. Le tableau moderne range les éléments par $Z$ croissant.
\end{savaistu}

\begin{experience}[Reconstituer la démarche de Mendeleïev (activité documentaire)]
\textbf{Matériel :} cartes « éléments » (symbole, masse atomique, propriétés : métal/non-métal, valence, aspect) pour les 20 premiers éléments.
\textbf{Protocole :}
\begin{enumerate}
\item Classer les cartes par masse atomique croissante.
\item Les ranger en lignes de façon que les éléments de propriétés voisines (ex. métaux très réactifs avec l'eau : Li, Na, K) tombent dans la même colonne.
\item Observer les cases vides laissées et prédire les propriétés de l'élément manquant entre le silicium et l'étain (germanium).
\end{enumerate}
\textbf{Interprétation :} la régularité (périodicité) permet de prévoir l'existence et les propriétés d'éléments inconnus. C'est la puissance du modèle de Mendeleïev.
\end{experience}

\courssub{Structure du tableau : lignes (périodes) et colonnes (familles)}

\begin{propriete}[Organisation du tableau de classification]
Le tableau de Mendeleïev (ou tableau périodique) range tous les éléments connus par \textbf{numéro atomique $Z$ croissant}.
\begin{itemize}
\item Les \cle{lignes horizontales} sont appelées \cle{périodes}. Le tableau en compte \textbf{7}. Le numéro de la période indique le nombre de couches électroniques occupées par l'atome.
\item Les \cle{colonnes verticales} sont appelées \cle{familles} (ou groupes). Le tableau en compte \textbf{18}. Les éléments d'une même famille ont le même nombre d'électrons sur leur couche externe (pour les familles principales) et donc des \textbf{propriétés chimiques voisines}.
\end{itemize}
\end{propriete}

\begin{popfigure}
\begin{tikzpicture}[scale=0.55, every node/.style={font=\small}]
% Grille 18 colonnes x 7 lignes
\draw[popline, thin] (0,0) grid (18,-7);
% Étiquettes Périodes (lignes)
\foreach \y/\n in {0/1, -1/2, -2/3, -3/4, -4/5, -5/6, -6/7}
  \node[left=2mm, popdark, font=\bfseries] at (-0.5, \y-0.5) {Période \n};
% Étiquettes Familles (colonnes)
\foreach \x/\n in {0.5/1, 1.5/2, 2.5/3, 3.5/4, 4.5/5, 5.5/6, 6.5/7, 7.5/8, 8.5/9, 9.5/10, 10.5/11, 11.5/12, 12.5/13, 13.5/14, 14.5/15, 15.5/16, 16.5/17, 17.5/18}
  \node[above=1mm, popdark, font=\bfseries] at (\x, 0.3) {\n};
% Coloration familles principales
\fill[popblue!20] (0,-1) rectangle (1,-6); % Col 1 (Alcalins) sans H
\fill[poporange!20] (1,-1) rectangle (2,-6); % Col 2 (Alcalino-terreux)
\fill[popgreen!20] (16,-1) rectangle (17,-6); % Col 17 (Halogènes)
\fill[poppurple!20] (17,-1) rectangle (18,-6); % Col 18 (Gaz rares)
% Bloc d (transition) - hachuré
\foreach \y in {3,4,5,6} {
  \fill[pattern=north east lines, pattern color=popmuted] (2,-\y) rectangle (12,-\y-1);
}
% Ligne de séparation Métaux / Non-métaux (escalier simplifié)
\draw[popdark, thick, decorate, decoration={zigzag, segment length=4, amplitude=1.5}] (12,-1) -- (13,-2) -- (14,-3) -- (15,-4) -- (16,-5) -- (17,-6);
% Légendes
\node[popblue, font=\bfseries] at (0.5, -6.7) {Alcalins};
\node[poporange, font=\bfseries] at (1.5, -6.7) {Alcalino-terreux};
\node[popgreen, font=\bfseries] at (16.5, -6.7) {Halogènes};
\node[poppurple, font=\bfseries] at (17.5, -6.7) {Gaz rares};
\node[popmuted, font=\bfseries, rotate=90] at (12.5, -3.5) {Métaux de transition};
\node[popdark, font=\bfseries] at (9, -8) {Métaux $\leftarrow$ \hspace{2cm} $\rightarrow$ Non-métaux};
% Hydrogène case spéciale
\fill[poppink!30] (0,0) rectangle (1,-1);
\node[popdark] at (0.5, -0.5) {\ce{H}};
\node[poppink, font=\tiny] at (0.5, -0.15) {Non-métal};
\end{tikzpicture}
\legende{Structure simplifiée du tableau de Mendeleïev : 7 périodes (lignes), 18 familles (colonnes). En couleur : familles principales (alcalins, alcalino-terreux, halogènes, gaz rares). Hachuré : métaux de transition. Ligne en zigzag : frontière métaux / non-métaux. L'hydrogène (H) est un cas à part.}
\end{popfigure}

\begin{methode}[Lire lignes, colonnes et numéro atomique]
Le tableau de Mendeleïev est organisé de manière rigoureuse :
\begin{enumerate}
\item Les éléments sont rangés par \cle{numéro atomique} $Z$ croissant, de gauche à droite puis de haut en bas.
\item Les \cle{lignes horizontales} s'appellent des \cle{périodes} : le tableau en compte 7.
\item Les \cle{colonnes verticales} s'appellent des \cle{familles} (ou groupes) : le tableau en compte 18.
\item Les éléments d'une même colonne ont des \cle{propriétés chimiques voisines}.
\end{enumerate}
\textbf{Réflexe :} pour situer un élément, donner toujours deux informations : sa période (ligne) et sa famille (colonne).
\end{methode}

\begin{exoresolu}[Situer des éléments dans le tableau]
\textbf{Énoncé.} Le sodium \ce{Na} ($Z = 11$) se trouve à la 3\up{e} période et dans la 1\up{re} colonne. Le chlore \ce{Cl} ($Z = 17$) se trouve à la 3\up{e} période, 17\up{e} colonne.
\begin{enumerate}
\item Que peut-on dire des propriétés chimiques du sodium et du potassium \ce{K} ($Z = 19$, 4\up{e} période, 1\up{re} colonne) ?
\item Le sodium et le chlore appartiennent-ils à la même famille ? Justifier.
\end{enumerate}
\tcblower
\textbf{Solution.}
\begin{enumerate}
\item Le sodium \ce{Na} et le potassium \ce{K} sont situés dans la \cle{même colonne} (1\up{re} colonne). Or les éléments d'une même colonne ont des propriétés chimiques voisines. Donc le sodium et le potassium ont des \textbf{propriétés chimiques semblables} : ce sont deux métaux très réactifs qui réagissent violemment avec l'eau.
\item Le sodium est dans la 1\up{re} colonne et le chlore dans la 17\up{e} colonne : ils n'appartiennent \textbf{pas à la même famille}. Ils sont seulement dans la même \cle{période} (la 3\up{e} ligne). Leurs propriétés sont très différentes : le sodium est un métal, le chlore est un non-métal gazeux.
\end{enumerate}
\textbf{Conclusion :} c'est la colonne (famille) qui renseigne sur la ressemblance des propriétés chimiques, pas la ligne.
\end{exoresolu}

\courssub{Quelques familles remarquables}

\begin{definition}[Familles principales du tableau (éléments représentatifs)]
On retient quatre familles majeures aux propriétés très marquées :
\begin{center}
\begin{tabularx}{\linewidth}{|l|X|X|X|}\hline
\rowcolor{popblueL} Famille & Colonne & Éléments types & Propriétés générales \\ \hline
\cle{Alcalins} & 1 (sauf H) & \ce{Li}, \ce{Na}, \ce{K}, \ce{Rb}, \ce{Cs} & Métaux mous, brillants, \textbf{très réactifs} avec l'eau (dégagent \ce{H2} et forment une base). \\ \hline
\cle{Alcalino-terreux} & 2 & \ce{Be}, \ce{Mg}, \ce{Ca}, \ce{Sr}, \ce{Ba} & Métaux gris, moins réactifs que les alcalins, présents dans les roches (calcaire \ce{CaCO3}, dolomie). \\ \hline
\cle{Halogènes} & 17 & \ce{F}, \ce{Cl}, \ce{Br}, \ce{I}, \ce{At} & Non-métaux \textbf{très réactifs}, forment des sels avec les métaux (ex. \ce{NaCl}). Existent sous forme diatomique (\ce{Cl2}, \ce{Br2}). \\ \hline
\cle{Gaz rares} (Gaz nobles) & 18 & \ce{He}, \ce{Ne}, \ce{Ar}, \ce{Kr}, \ce{Xe} & Gaz monoatomiques, \textbf{inertes} (peu ou pas de réactions chimiques), couche électronique saturée. \\ \hline
\end{tabularx}
\end{center}
\end{definition}

\begin{methode}[Identifier les grandes familles du tableau]
Pour reconnaître une famille remarquable :
\begin{enumerate}
\item \cle{Alcalins} : 1\up{re} colonne (sauf \ce{H}) : \ce{Li}, \ce{Na}, \ce{K}... Métaux mous, réagissent violemment avec l'eau.
\item \cle{Alcalino-terreux} : 2\up{e} colonne : \ce{Mg}, \ce{Ca}... Métaux présents dans les roches (calcaire).
\item \cle{Halogènes} : 17\up{e} colonne : \ce{F}, \ce{Cl}, \ce{Br}, \ce{I}. Non-métaux très réactifs (désinfectants, sel de cuisine).
\item \cle{Gaz rares} (ou gaz nobles) : 18\up{e} colonne : \ce{He}, \ce{Ne}, \ce{Ar}. Gaz \cle{inertes} : ils ne réagissent presque jamais.
\end{enumerate}
\textbf{Réflexe :} mémoriser la position de chaque famille (numéro de colonne) et un exemple concret d'usage au Cameroun (ex. chlore eau SONARA/ENEO, néon enseignes Douala/Yaoundé, calcium ciment CIMENCAM, aluminium marmites).
\end{methode}

\begin{exoresolu}[Familles et usages quotidiens]
\textbf{Énoncé.} À Douala, on utilise le néon pour les enseignes lumineuses, le chlore pour traiter l'eau de la piscine et le calcium est abondant dans le calcaire. Pour chaque élément souligné (néon, chlore, calcium), donner sa famille et justifier son usage par ses propriétés.
\tcblower
\textbf{Solution.}
\begin{enumerate}
\item \textbf{Néon \ce{Ne}} : il appartient à la famille des \cle{gaz rares} (18\up{e} colonne). C'est un gaz inerte qui ne réagit pas avec les autres corps ; soumis à une décharge électrique, il émet une lumière rouge-orange, d'où son usage dans les enseignes lumineuses.
\item \textbf{Chlore \ce{Cl}} : il appartient à la famille des \cle{halogènes} (17\up{e} colonne). C'est un non-métal très réactif, capable de détruire les microbes : on l'utilise donc pour désinfecter l'eau des piscines et l'eau de boisson (traitement SONARA/ENEO).
\item \textbf{Calcium \ce{Ca}} : il appartient à la famille des \cle{alcalino-terreux} (2\up{e} colonne). Il est très abondant dans l'écorce terrestre, notamment dans le calcaire \ce{CaCO3}, matériau de base pour le ciment (CIMENCAM) et la construction.
\end{enumerate}
\textbf{Conclusion :} la famille d'un élément permet de prévoir ses propriétés et donc ses applications.
\end{exoresolu}

\courssub{Éléments usuels et applications dans la vie courante}

\begin{methode}[Utiliser le numéro atomique $Z$ pour la composition de l'atome]
Le numéro atomique $Z$ d'un élément est le \cle{nombre de protons} de son noyau. Pour un atome (électriquement neutre) :
\begin{enumerate}
\item Nombre de protons $= Z$.
\item Nombre d'électrons $= Z$ (car l'atome est neutre : autant de charges $+$ que de charges $-$).
\item On répartit les électrons sur les couches (K, L, M, N...) : 2 max sur K, 8 max sur L, 8 max sur M (pour les 20 premiers éléments), etc.
\end{enumerate}
\textbf{Réflexe :} le numéro de la période correspond au nombre de couches électroniques occupées ; le numéro de la colonne (pour les familles 1, 2, 13 à 18) est lié au nombre d'électrons de la couche externe.
\end{methode}

\begin{popfigure}
\begin{tikzpicture}[scale=1.2]
% Atome de Magnésium Z=12
% Noyau
\draw[popdark, thick, fill=poporange!30] (0,0) circle (0.6);
\node[popdark, font=\bfseries, align=center] at (0,0){Noyau\\$12p^+$, $12n^0$};
% Couche K
\draw[popblue, thick] (0,0) circle (1.2);
\foreach \angle in {90, 210, 330} \fill[popblue] (\angle:1.2) circle (0.08);
\node[popblue] at (0, 1.5) {Couche K : 2 $e^-$};
% Couche L
\draw[popgreen, thick] (0,0) circle (2.0);
\foreach \angle in {18, 90, 162, 234, 306, 30, 150, 270} \fill[popgreen] (\angle:2.0) circle (0.08);
\node[popgreen] at (0, 2.4) {Couche L : 8 $e^-$};
% Couche M
\draw[poppurple, thick] (0,0) circle (2.8);
\foreach \angle in {90, 270} \fill[poppurple] (\angle:2.8) circle (0.08);
\node[poppurple] at (0, 3.2) {Couche M : 2 $e^-$};
% Légende
\node[align=center, font=\small] at (4.5, 0) {\textbf{Magnésium \ce{Mg} ($Z=12$)}\\12 protons ($p^+$)\\12 neutrons ($n^0$)\\12 électrons ($e^-$)\\Répartition : \textbf{2, 8, 2}\\Période 3 (3 couches)\\Famille 2 (2 $e^-$ externes)};
\end{tikzpicture}
\legende{Modèle de Bohr simplifié de l'atome de magnésium (Z=12). La répartition électronique 2, 8, 2 place Mg en Période 3 (3 couches) et Famille 2 (2 électrons externes).}
\end{popfigure}

\begin{exoresolu}[Composition de l'atome de magnésium]
\textbf{Énoncé.} Le magnésium \ce{Mg} a pour numéro atomique $Z = 12$.
\begin{enumerate}
\item Donner le nombre de protons et d'électrons de l'atome de magnésium.
\item Écrire la répartition électronique de cet atome.
\item En déduire sa période dans le tableau.
\end{enumerate}
\tcblower
\textbf{Solution.}
\begin{enumerate}
\item L'atome de magnésium possède $Z = 12$ \textbf{protons}. Comme l'atome est électriquement neutre, il possède aussi \textbf{12 électrons}.
\item Répartition des électrons : la 1\up{re} couche (K) reçoit 2 électrons (maximum), la 2\up{e} couche (L) reçoit 8 électrons (maximum), il reste $12 - 2 - 8 = 2$ électrons sur la 3\up{e} couche (M). On écrit : $(K)^2\,(L)^8\,(M)^2$ ou simplement \textbf{2, 8, 2}.
\item L'atome occupe \textbf{3 couches électroniques} : le magnésium se trouve donc à la \textbf{3\up{e} période} du tableau. Avec 2 électrons sur la couche externe, il est dans la 2\up{e} colonne : c'est bien un alcalino-terreux.
\end{enumerate}
\textbf{Conclusion :} \ce{Mg} : 12 protons, 12 électrons répartis 2, 8, 2 ; il appartient à la 3\up{e} période et à la famille des alcalino-terreux.
\end{exoresolu}

\begin{methode}[Classer un élément : métal ou non-métal]
\begin{enumerate}
\item Repérer la position de l'élément dans le tableau.
\item Les \cle{métaux} occupent la partie \cle{gauche et centrale} du tableau (alcalins, alcalino-terreux, métaux de transition comme \ce{Fe}, \ce{Cu}, \ce{Zn}).
\item Les \cle{non-métaux} occupent la partie \cle{droite} du tableau (\ce{C}, \ce{N}, \ce{O}, \ce{P}, \ce{S}, halogènes, gaz rares).
\item La frontière est une ligne en escalier (zigzag) entre les colonnes 13 et 17 (passant par \ce{B}, \ce{Si}, \ce{Ge}, \ce{As}, \ce{Sb}, \ce{Te}, \ce{Po}).
\item Vérifier avec les propriétés : les métaux conduisent le courant électrique et la chaleur, sont brillants, malléables et ductiles.
\end{enumerate}
\textbf{Réflexe :} l'hydrogène \ce{H}, bien que placé en 1\up{re} colonne, est un \textbf{non-métal} : c'est une exception à connaître. Les métalloïdes (\ce{B}, \ce{Si}, \ce{Ge}, \ce{As}, \ce{Sb}, \ce{Te}) ont des propriétés intermédiaires.
\end{methode}

\begin{exoresolu}[Métaux et non-métaux autour de nous]
\textbf{Énoncé.} Classer les éléments suivants en métaux et non-métaux, puis donner une application courante de chacun : \ce{Cu}, \ce{O}, \ce{Al}, \ce{Cl}, \ce{Fe}, \ce{C}.
\tcblower
\textbf{Solution.}
\begin{enumerate}
\item \ce{Cu} (cuivre) : situé dans la partie centrale du tableau (métaux de transition) $\Rightarrow$ \textbf{métal}. Application : fils électriques des installations ENEO (excellent conducteur).
\item \ce{O} (oxygène) : partie droite du tableau $\Rightarrow$ \textbf{non-métal}. Application : gaz de l'air indispensable à la respiration et aux combustions (cuisine au bois/charbon).
\item \ce{Al} (aluminium) : \textbf{métal} (famille 13). Application : marmites, casseroles et tôles de toiture (légèreté, conduction chaleur).
\item \ce{Cl} (chlore) : halogène, partie droite $\Rightarrow$ \textbf{non-métal}. Application : désinfection de l'eau (eau de robinet, piscines).
\item \ce{Fe} (fer) : métal de transition $\Rightarrow$ \textbf{métal}. Application : fers à béton pour la construction des maisons, carrosseries de moto-taxi.
\item \ce{C} (carbone) : partie droite $\Rightarrow$ \textbf{non-métal}. Application : charbon de bois pour la cuisine, graphite dans les piles.
\end{enumerate}
\textbf{Conclusion :} métaux : \ce{Cu}, \ce{Al}, \ce{Fe} ; non-métaux : \ce{O}, \ce{Cl}, \ce{C}.
\end{exoresolu}

\begin{methode}[Réussir une question de type BEPC sur la classification]
\begin{enumerate}
\item \textbf{Lire} l'énoncé en entier et repérer les données : symboles, numéros atomiques, positions.
\item \textbf{Situer} chaque élément : période (ligne) et famille (colonne).
\item \textbf{Comparer} : même colonne $\Rightarrow$ propriétés voisines ; même ligne $\Rightarrow$ même période seulement.
\item \textbf{Rédiger} la réponse en phrases complètes, avec le vocabulaire exact : \emph{élément}, \emph{symbole}, \emph{numéro atomique}, \emph{période}, \emph{famille}.
\item \textbf{Conclure} par une phrase qui répond précisément à la question posée.
\end{enumerate}
\textbf{Réflexe :} au BEPC, toute affirmation doit être justifiée par une donnée du tableau (numéro de ligne ou de colonne).
\end{methode}

\begin{exoresolu}[Sujet type BEPC]
\textbf{Énoncé.} On donne les numéros atomiques suivants : lithium \ce{Li} ($Z = 3$), béryllium \ce{Be} ($Z = 4$), sodium \ce{Na} ($Z = 11$), magnésium \ce{Mg} ($Z = 12$).
\begin{enumerate}
\item Donner la répartition électronique de chacun de ces atomes.
\item Placer ces éléments dans le tableau (période et colonne).
\item Quels éléments ont des propriétés chimiques voisines ? Justifier.
\end{enumerate}
\tcblower
\textbf{Solution.}
\begin{enumerate}
\item Répartitions électroniques (règle 2, 8, 8...) :
\begin{itemize}
\item \ce{Li} ($Z = 3$) : 3 électrons $\Rightarrow$ \textbf{2, 1} ;
\item \ce{Be} ($Z = 4$) : 4 électrons $\Rightarrow$ \textbf{2, 2} ;
\item \ce{Na} ($Z = 11$) : 11 électrons $\Rightarrow$ \textbf{2, 8, 1} ;
\item \ce{Mg} ($Z = 12$) : 12 électrons $\Rightarrow$ \textbf{2, 8, 2}.
\end{itemize}
\item Placement dans le tableau (nombre de couches = période ; électrons externes = colonne pour familles 1 et 2) :
\begin{center}
\begin{tabularx}{\linewidth}{|l|X|X|X|}\hline
\rowcolor{popblueL} Élément & Répartition & Période & Colonne (famille) \\ \hline
\ce{Li} & 2, 1 & 2\up{e} & 1\up{re} (alcalin) \\ \hline
\ce{Be} & 2, 2 & 2\up{e} & 2\up{e} (alcalino-terreux) \\ \hline
\ce{Na} & 2, 8, 1 & 3\up{e} & 1\up{re} (alcalin) \\ \hline
\ce{Mg} & 2, 8, 2 & 3\up{e} & 2\up{e} (alcalino-terreux) \\ \hline
\end{tabularx}
\end{center}
\item \ce{Li} et \ce{Na} sont dans la même colonne (1\up{re}) : ce sont des \textbf{alcalins}, aux propriétés voisines (métaux réagissant violemment avec l'eau). De même, \ce{Be} et \ce{Mg} (2\up{e} colonne) sont des alcalino-terreux aux propriétés voisines.
\end{enumerate}
\textbf{Conclusion :} les éléments de même famille (même colonne, même nombre d'électrons externes) ont des propriétés chimiques semblables : c'est le principe fondateur de la classification de Mendeleïev.
\end{exoresolu}

\begin{attention}[Les 5 erreurs qui coûtent des points au BEPC]
\begin{enumerate}
\item \textbf{Mal écrire les symboles} : écrire \ce{NA}, \ce{CL} ou \ce{FE} est faux. Un symbole s'écrit avec \textbf{une majuscule} suivie éventuellement d'\textbf{une minuscule} : \ce{Na}, \ce{Cl}, \ce{Fe}.
\item \textbf{Confondre période et famille} : la période est la \textbf{ligne horizontale}, la famille est la \textbf{colonne verticale}. Dire que deux éléments de la même ligne ont des propriétés voisines est faux : seule la colonne compte pour la ressemblance chimique.
\item \textbf{Confondre numéro atomique et nombre de masse} : le numéro atomique $Z$ est le nombre de \textbf{protons} (et d'électrons dans l'atome neutre), pas le nombre total de nucléons ($A = Z + N$). Ne pas additionner protons et neutrons pour trouver $Z$.
\item \textbf{Oublier que l'hydrogène est un non-métal} : bien que placé en tête de la 1\up{re} colonne, \ce{H} n'est \textbf{pas} un alcalin. Citer \ce{H} comme exemple d'alcalin est une erreur classique.
\item \textbf{Répondre sans justifier par le tableau} : affirmer « le sodium et le potassium sont semblables » sans préciser qu'ils sont dans la \textbf{même colonne} (1\up{re}) fait perdre le point de justification. Toujours citer la période et la famille.
\end{enumerate}
\end{attention}

\newpage

\coursec{Exercices}

\courssub{Je vérifie mes connaissances}

\begin{exercice}[QCM : le tableau périodique]
Pour chaque question, entoure la (ou les) bonne(s) réponse(s).
\begin{enumerate}
\item Dans le tableau de classification périodique, les éléments sont rangés par ordre croissant de :
\begin{enumerate}
\item masse ;
\item numéro atomique $Z$ ;
\item date de découverte.
\end{enumerate}
\item Une ligne horizontale du tableau s'appelle :
\begin{enumerate}
\item une famille ;
\item une période ;
\item un groupe.
\end{enumerate}
\item Le symbole chimique du fer est :
\begin{enumerate}
\item \ce{FE} ;
\item \ce{Fe} ;
\item \ce{F}.
\end{enumerate}
\item Les éléments d'une même colonne :
\begin{enumerate}
\item ont le même numéro atomique ;
\item ont des propriétés chimiques voisines ;
\item appartiennent à la même période.
\end{enumerate}
\end{enumerate}
\end{exercice}

\begin{exercice}[Vrai ou faux ? Justifie]
Réponds par \textbf{vrai} ou \textbf{faux}, puis justifie ta réponse en une ou deux phrases.
\begin{enumerate}
\item Le symbole du cobalt est \ce{CO}.
\item Le chlore \ce{Cl} et le fluor \ce{F} appartiennent à la même famille.
\item Le tableau de Mendeleïev ne contient que des éléments métalliques.
\item Le numéro atomique d'un élément est égal au nombre de protons de son atome.
\end{enumerate}
\end{exercice}

\begin{exercice}[Définitions]
Définis en une phrase chacun des termes suivants, en donnant un exemple :
\begin{enumerate}
\item élément chimique ;
\item période ;
\item famille (ou colonne) ;
\item symbole chimique.
\end{enumerate}
\end{exercice}

\begin{exercice}[Symboles et numéros atomiques]
Recopie et complète le tableau suivant :
\begin{center}
\begin{tabularx}{\linewidth}{|l|X|X|}\hline
\rowcolor{popblueL} \textbf{Nom de l'élément} & \textbf{Symbole} & \textbf{Numéro atomique $Z$} \\ \hline
Oxygène & & 8 \\ \hline
 & \ce{Na} & 11 \\ \hline
Calcium & \ce{Ca} & \\ \hline
 & \ce{Cl} & 17 \\ \hline
\end{tabularx}
\end{center}
\end{exercice}

\courssub{Je m'entraîne}

\begin{exercice}[Lecture d'un extrait du tableau]
On donne l'extrait simplifié du tableau périodique ci-dessous (les numéros atomiques sont indiqués entre parenthèses) :
\begin{center}
\begin{tabularx}{\linewidth}{|l|X|X|X|X|}\hline
\rowcolor{popblueL} \textbf{Période} & \textbf{Colonne 1} & \textbf{Colonne 2} & \textbf{Colonne 16} & \textbf{Colonne 17} \\ \hline
2 & \ce{Li} (3) & \ce{Be} (4) & \ce{O} (8) & \ce{F} (9) \\ \hline
3 & \ce{Na} (11) & \ce{Mg} (12) & \ce{S} (16) & \ce{Cl} (17) \\ \hline
4 & \ce{K} (19) & \ce{Ca} (20) & \ce{Se} (34) & \ce{Br} (35) \\ \hline
\end{tabularx}
\end{center}
\begin{enumerate}
\item Donne le symbole, le numéro atomique, la période et la colonne de chacun des éléments suivants : sodium, chlore, calcium, oxygène, potassium.
\item Quels éléments de cet extrait appartiennent à la même famille que le sodium ?
\item Quels éléments appartiennent à la même période que le chlore ?
\end{enumerate}
\end{exercice}

\begin{exercice}[Éléments et usages quotidiens]
Associe chaque élément de la colonne de gauche à son usage quotidien dans la colonne de droite, puis justifie chaque association par une propriété de l'élément.
\begin{center}
\begin{tabularx}{\linewidth}{|l|X|}\hline
\rowcolor{popblueL} \textbf{Élément} & \textbf{Usage} \\ \hline
Aluminium \ce{Al} & Fabrication des marmites et marmites de cuisine \\ \hline
Cuivre \ce{Cu} & Câbles électriques des installations ENEO \\ \hline
Chlore \ce{Cl} & Désinfection de l'eau de boisson (eau de Javel) \\ \hline
Calcium \ce{Ca} & Solidité des os et des dents \\ \hline
Fer \ce{Fe} & Tiges à béton des constructions \\ \hline
\end{tabularx}
\end{center}
\end{exercice}

\begin{exercice}[L'erreur de Junior]
En recopiant son cours, Junior écrit : « \ce{CO} est le symbole du cobalt ».
\begin{enumerate}
\item Explique pourquoi cette écriture est fausse et ce que représente réellement \ce{CO}.
\item Écris le symbole correct du cobalt.
\item Énonce la règle d'écriture des symboles chimiques (nombre de lettres, majuscules et minuscules).
\item Applique cette règle pour corriger les écritures suivantes : \ce{CL}, \ce{cA}, \ce{MG}.
\end{enumerate}
\end{exercice}

\begin{exercice}[Sodium et potassium : même famille]
Le sodium \ce{Na} ($Z = 11$) et le potassium \ce{K} ($Z = 19$) sont situés dans la première colonne du tableau périodique.
\begin{enumerate}
\item Quel est le nom de la famille d'éléments à laquelle ils appartiennent ?
\item Que peut-on prévoir de leurs propriétés chimiques ? Justifie.
\item Le sodium réagit violemment avec l'eau. Que peut-on prévoir pour le potassium ?
\item Cite une précaution de manipulation et de conservation du sodium au laboratoire.
\end{enumerate}
\end{exercice}

\courssub{Je me prépare au Bac}

\begin{exercice}[Type BEPC : exploitation du tableau périodique]
\textit{(D'après un sujet de BEPC, session 2019)}

On considère les cinq éléments chimiques suivants : sodium \ce{Na} ($Z = 11$), chlore \ce{Cl} ($Z = 17$), calcium \ce{Ca} ($Z = 20$), fer \ce{Fe} ($Z = 26$) et oxygène \ce{O} ($Z = 8$). Leur position dans le tableau périodique est rappelée : \ce{O} et \ce{S} sont dans la même colonne ; \ce{Na} et \ce{K} sont dans la même colonne ; \ce{Cl} et \ce{F} sont dans la même colonne ; \ce{Ca} est dans la colonne qui suit celle de \ce{K} ; \ce{Fe} est un métal de transition de la 4\up{e} période.
\begin{enumerate}
\item Rappelle ce que représente le numéro atomique $Z$ d'un élément.
\item Pour chacun des cinq éléments, indique le numéro atomique, la période et la famille (colonne) auxquelles il appartient. Présente les résultats dans un tableau.
\item Parmi ces cinq éléments, lesquels sont des métaux ? Lesquels sont des non-métaux ?
\item Le chlore est utilisé à Yaoundé pour traiter l'eau de forage avant consommation. À quelle famille appartient-il ? Cite un autre élément de cette famille.
\item Le fer rouille au contact de l'air humide. Propose une méthode simple utilisée au Cameroun pour protéger les tiges à béton ou les tôles en fer.
\end{enumerate}
\end{exercice}

\begin{exercice}[Type BEPC : la prédiction de Mendeleïev]
\textit{Document.} En 1871, Dmitri Mendeleïev laisse des cases vides dans son tableau et prédit l'existence d'éléments inconnus. Pour l'« eka-aluminium », placé sous l'aluminium, il prédit : masse atomique voisine de 68, densité voisine de 5,9, métal fusible à basse température. En 1875, le Français Paul-Émile Lecoq de Boisbaudran découvre le gallium \ce{Ga} : masse atomique 69,7, densité mesurée 5,94, température de fusion 29,8 \degres C (il fond dans la main !).
\begin{enumerate}
\item Qui est Mendeleïev et quelle est sa contribution majeure à la chimie ?
\item Sur quel critère Mendeleïev rangeait-il les éléments dans son tableau ? Quel critère est utilisé dans le tableau moderne ?
\item Compare les propriétés prédites de l'eka-aluminium et les propriétés mesurées du gallium (masse atomique, densité). Calcule l'écart relatif entre la densité prédite et la densité mesurée, en pourcentage : $\dfrac{5,94 - 5,9}{5,94} \times 100$.
\item Explique pourquoi la découverte du gallium a constitué une validation spectaculaire de la classification périodique.
\item Le gallium \ce{Ga} ($Z = 31$) est placé sous l'aluminium \ce{Al} ($Z = 13$). À quelle famille appartiennent ces deux éléments ? Que peut-on dire de leurs propriétés chimiques ?
\end{enumerate}
\end{exercice}

\begin{exercice}[Type BEPC : sécurité domestique et famille des halogènes]
\textit{Situation.} Maman Aïcha veut nettoyer les toilettes de la maison. Elle verse de l'eau de Javel (qui contient des ions hypochlorite \ce{ClO-}) puis, sans attendre, un détartrant acide (qui contient des ions \ce{H+}). Un dégagement gazeux verdâtre et irritant se produit immédiatement, et Maman Aïcha se met à tousser.

L'équation de la réaction est : \ce{ClO- + Cl- + 2H+ -> Cl2 + H2O}
\begin{enumerate}
\item Nomme le gaz dégagé lors de ce mélange et donne sa formule chimique.
\item À quelle famille du tableau périodique appartient l'élément chlore ? Cite deux autres éléments de cette famille.
\item Donne le numéro atomique du chlore ($Z = 17$) et précise sa période dans le tableau.
\item Explique, en t'appuyant sur l'équation de la réaction, pourquoi il ne faut jamais mélanger l'eau de Javel avec un détartrant acide.
\item Rédige une consigne de sécurité en deux phrases que tu afficherais dans la cuisine familiale concernant l'utilisation et le rangement des produits d'entretien.
\end{enumerate}
\end{exercice}

\newpage

\coursec{Corrigés des exercices}

\begin{corrige}[Exercice 1 — QCM : le tableau périodique]
\begin{enumerate}
\item \textbf{Bonne réponse : b.} Dans le tableau moderne, les éléments sont classés par ordre croissant de \cle{numéro atomique} $Z$ (nombre de protons). Mendeleïev utilisait les masses atomiques, mais le critère actuel est $Z$.
\item \textbf{Bonne réponse : b.} Une ligne horizontale s'appelle une \cle{période}. Une colonne verticale s'appelle une famille (ou groupe).
\item \textbf{Bonne réponse : b.} Le symbole du fer est \ce{Fe} (du latin \emph{ferrum}) : une majuscule suivie d'une minuscule. \ce{FE} est incorrect (deux majuscules) et \ce{F} est le symbole du fluor.
\item \textbf{Bonne réponse : b.} Les éléments d'une même colonne ont des \cle{propriétés chimiques voisines} (même nombre d'électrons sur leur couche externe). Ils n'ont pas le même numéro atomique et n'appartiennent pas à la même période.
\end{enumerate}
\end{corrige}

\begin{corrige}[Exercice 2 — Vrai ou faux ? Justifie]
\begin{enumerate}
\item \textbf{Faux.} Le symbole du cobalt est \ce{Co} (une majuscule puis une minuscule). \ce{CO} représente la molécule de monoxyde de carbone, formée d'un atome de carbone \ce{C} et d'un atome d'oxygène \ce{O}.
\item \textbf{Vrai.} Le chlore \ce{Cl} et le fluor \ce{F} sont situés dans la même colonne (la 17\textsuperscript{e} colonne) : ils appartiennent à la famille des \cle{halogènes}.
\item \textbf{Faux.} Le tableau contient des métaux (fer, sodium, cuivre...), mais aussi des non-métaux (oxygène, chlore, carbone...) et des métalloïdes (silicium...).
\item \textbf{Vrai.} Par définition, le numéro atomique $Z$ d'un élément est le nombre de \cle{protons} contenus dans le noyau de son atome. C'est aussi le nombre d'électrons de l'atome, qui est électriquement neutre.
\end{enumerate}
\end{corrige}

\begin{corrige}[Exercice 3 — Définitions]
\begin{enumerate}
\item \textbf{Élément chimique :} ensemble d'atomes (ou d'ions) possédant le même nombre de protons, c'est-à-dire le même numéro atomique $Z$. \emph{Exemple :} l'élément oxygène ($Z = 8$).
\item \textbf{Période :} ligne horizontale du tableau périodique ; les éléments d'une même période ont le même nombre de couches électroniques. \emph{Exemple :} le sodium \ce{Na} et le chlore \ce{Cl} appartiennent à la 3\textsuperscript{e} période.
\item \textbf{Famille (ou colonne) :} colonne verticale du tableau périodique ; les éléments d'une même famille ont des propriétés chimiques voisines. \emph{Exemple :} la famille des alcalins (\ce{Li}, \ce{Na}, \ce{K}...).
\item \textbf{Symbole chimique :} écriture abrégée d'un élément, formée d'une ou deux lettres (la première en majuscule, la deuxième en minuscule). \emph{Exemple :} \ce{Fe} pour le fer.
\end{enumerate}
\end{corrige}

\begin{corrige}[Exercice 4 — Symboles et numéros atomiques]
\begin{center}
\begin{tabularx}{\linewidth}{|l|X|X|}\hline
\rowcolor{popblueL} \textbf{Nom de l'élément} & \textbf{Symbole} & \textbf{Numéro atomique $Z$} \\ \hline
Oxygène & \ce{O} & 8 \\ \hline
Sodium & \ce{Na} & 11 \\ \hline
Calcium & \ce{Ca} & 20 \\ \hline
Chlore & \ce{Cl} & 17 \\ \hline
\end{tabularx}
\end{center}
\emph{Justifications :} le symbole de l'oxygène est \ce{O} (première lettre du nom). \ce{Na} vient du latin \emph{natrium} : c'est le sodium. Le calcium a pour numéro atomique $Z = 20$ (lecture du tableau). \ce{Cl} est le symbole du chlore.
\end{corrige}

\begin{corrige}[Exercice 5 — Lecture d'un extrait du tableau]
\begin{enumerate}
\item Lecture directe de l'extrait :
\begin{center}
\begin{tabularx}{\linewidth}{|l|X|X|X|}\hline
\rowcolor{popblueL} \textbf{Élément} & \textbf{Symbole ($Z$)} & \textbf{Période} & \textbf{Colonne} \\ \hline
Sodium & \ce{Na} (11) & 3 & 1 \\ \hline
Chlore & \ce{Cl} (17) & 3 & 17 \\ \hline
Calcium & \ce{Ca} (20) & 4 & 2 \\ \hline
Oxygène & \ce{O} (8) & 2 & 16 \\ \hline
Potassium & \ce{K} (19) & 4 & 1 \\ \hline
\end{tabularx}
\end{center}
\item Les éléments de la même famille que le sodium sont ceux de la colonne 1 : le \cle{lithium} \ce{Li} et le \cle{potassium} \ce{K}. Ce sont les alcalins.
\item Les éléments de la même période que le chlore (période 3) sont : \ce{Na}, \ce{Mg}, \ce{Al}, \ce{Si}, \ce{P}, \ce{S}, \ce{Cl}, \ce{Ar}.
\end{enumerate}
\end{corrige}

\begin{corrige}[Exercice 6 — Éléments et usages quotidiens]
\begin{itemize}
\item \textbf{Aluminium \ce{Al} — marmites de cuisine :} l'aluminium est un métal léger, bon conducteur de la chaleur et peu coûteux, ce qui convient à la cuisson des aliments.
\item \textbf{Cuivre \ce{Cu} — câbles électriques ENEO :} le cuivre est un excellent conducteur du courant électrique et il est ductile (on peut l'étirer en fils).
\item \textbf{Chlore \ce{Cl} — désinfection de l'eau (eau de Javel) :} le chlore est un puissant désinfectant qui tue les microbes présents dans l'eau de boisson.
\item \textbf{Calcium \ce{Ca} — solidité des os et des dents :} le calcium entre dans la constitution du tissu osseux et de l'émail dentaire, qu'il rend durs et résistants.
\item \textbf{Fer \ce{Fe} — tiges à béton :} le fer (sous forme d'acier) est un métal très résistant mécaniquement, idéal pour renforcer le béton des constructions.
\end{itemize}
\end{corrige}

\begin{corrige}[Exercice 7 — L'erreur de Junior]
\begin{enumerate}
\item \ce{CO} n'est pas le symbole d'un élément : c'est la formule d'une \cle{molécule}, le monoxyde de carbone, formée d'un atome de carbone \ce{C} et d'un atome d'oxygène \ce{O}. Un symbole d'élément ne comporte jamais deux majuscules.
\item Le symbole correct du cobalt est \ce{Co} (majuscule \ce{C}, minuscule \ce{o}).
\item \textbf{Règle d'écriture :} le symbole d'un élément chimique comporte une ou deux lettres ; la première lettre est toujours une \cle{majuscule} et la deuxième, si elle existe, est toujours une \cle{minuscule}.
\item Corrections : \ce{CL} $\rightarrow$ \ce{Cl} (chlore) ; \ce{cA} $\rightarrow$ \ce{Ca} (calcium) ; \ce{MG} $\rightarrow$ \ce{Mg} (magnésium).
\end{enumerate}
\end{corrige}

\begin{corrige}[Exercice 8 — Sodium et potassium : même famille]
\begin{enumerate}
\item Le sodium et le potassium appartiennent à la famille des \cle{alcalins} (1\textsuperscript{re} colonne du tableau).
\item Les éléments d'une même famille ont des propriétés chimiques voisines : on peut donc prévoir que le potassium se comporte chimiquement comme le sodium (métaux mous, très réactifs, réagissant avec l'eau).
\item Le potassium réagit aussi avec l'eau, et même \textbf{plus violemment} que le sodium : la réactivité des alcalins augmente quand on descend dans la colonne.
\item Précautions : le sodium se conserve \textbf{dans l'huile (ou le pétrole)} pour l'isoler de l'air et de l'humidité ; on le manipule avec une pince (jamais à mains nues), en portant des lunettes de protection, et loin de toute eau.
\end{enumerate}
\end{corrige}

\begin{corrige}[Exercice 9 — Type BEPC : exploitation du tableau périodique]
\begin{enumerate}
\item Le numéro atomique $Z$ d'un élément est le nombre de \cle{protons} contenus dans le noyau de son atome (égal au nombre d'électrons de l'atome neutre).
\item Tableau des positions :
\begin{center}
\begin{tabularx}{\linewidth}{|l|X|X|X|}\hline
\rowcolor{popblueL} \textbf{Élément} & \textbf{$Z$} & \textbf{Période} & \textbf{Famille (colonne)} \\ \hline
Oxygène \ce{O} & 8 & 2 & Colonne de \ce{S} : chalcogènes (16\textsuperscript{e} colonne) \\ \hline
Sodium \ce{Na} & 11 & 3 & Colonne de \ce{K} : alcalins (1\textsuperscript{re} colonne) \\ \hline
Chlore \ce{Cl} & 17 & 3 & Colonne de \ce{F} : halogènes (17\textsuperscript{e} colonne) \\ \hline
Calcium \ce{Ca} & 20 & 4 & Colonne suivant celle de \ce{K} : alcalino-terreux (2\textsuperscript{e} colonne) \\ \hline
Fer \ce{Fe} & 26 & 4 & Métaux de transition \\ \hline
\end{tabularx}
\end{center}
\emph{Justification des périodes :} $Z = 8$ et $Z = 11$ à $17$ remplissent les périodes 2 et 3 ; $Z = 20$ et $Z = 26$ appartiennent à la 4\textsuperscript{e} période.
\item \textbf{Métaux :} sodium \ce{Na}, calcium \ce{Ca}, fer \ce{Fe}. \textbf{Non-métaux :} chlore \ce{Cl}, oxygène \ce{O}.
\item Le chlore appartient à la famille des \cle{halogènes} (17\textsuperscript{e} colonne). Autre élément de cette famille : le fluor \ce{F} (ou le brome \ce{Br}, l'iode \ce{I}).
\item Pour protéger le fer de la rouille, on peut le \textbf{peindre} (peinture antirouille sur les tôles), le \textbf{graisser ou huiler}, ou le recouvrir d'une couche de zinc (\textbf{galvanisation}, utilisée pour les tôles « zinc » des toitures). Ces méthodes isolent le fer de l'air et de l'humidité.
\end{enumerate}
\textbf{Barème indicatif (sur 10) :} question 1 : 1 pt ; question 2 : 3 pts (0,5 pt par ligne correcte + cohérence du tableau) ; question 3 : 2 pts ; question 4 : 2 pts ; question 5 : 2 pts.

\textbf{Points de vigilance :} exiger la définition complète de $Z$ (protons du noyau, pas « nombre d'atomes ») ; ne pas confondre période (ligne) et famille (colonne) ; pour la question 5, la réponse doit mentionner une méthode concrète et le fait qu'elle empêche le contact avec l'air humide.
\end{corrige}

\begin{corrige}[Exercice 10 — Type BEPC : la prédiction de Mendeleïev]
\begin{enumerate}
\item Dmitri Mendeleïev est un chimiste russe qui, en 1869-1871, construisit le premier \cle{tableau de classification périodique} des éléments, en laissant des cases vides pour des éléments encore inconnus dont il prévit les propriétés.
\item Mendeleïev rangeait les éléments par ordre de \cle{masse atomique croissante} (en regroupant ceux de propriétés voisines). Dans le tableau moderne, le critère est le \cle{numéro atomique $Z$ croissant}.
\item Comparaison :
\begin{center}
\begin{tabularx}{\linewidth}{|l|X|X|}\hline
\rowcolor{popblueL} \textbf{Propriété} & \textbf{Eka-aluminium (prédit)} & \textbf{Gallium (mesuré)} \\ \hline
Masse atomique & $\approx 68$ & 69,7 \\ \hline
Densité & $\approx 5,9$ & 5,94 \\ \hline
Fusion & basse température & \SI{29,8}{\degreeCelsius} (fond dans la main) \\ \hline
\end{tabularx}
\end{center}
Écart relatif sur la densité :
\[ \frac{5,94 - 5,9}{5,94} \times 100 = \frac{0,04}{5,94} \times 100 \approx 0,67\,\% \]
L'écart est inférieur à $1\,\%$ : la prédiction était remarquablement précise.
\item La découverte du gallium, dont les propriétés mesurées correspondent presque exactement aux prédictions, prouve que la classification de Mendeleïev n'était pas un simple rangement mais un outil \textbf{prédictif} fondé sur une loi naturelle : c'est une validation expérimentale spectaculaire du tableau périodique.
\item \ce{Ga} et \ce{Al} sont dans la même colonne : ils appartiennent à la même \cle{famille} (colonne 13). Leurs propriétés chimiques sont donc \textbf{voisines} (même nombre d'électrons sur la couche externe).
\end{enumerate}
\textbf{Barème indicatif (sur 10) :} question 1 : 1,5 pt ; question 2 : 2 pts ; question 3 : 2,5 pts (tableau 1 pt, calcul 1,5 pt) ; question 4 : 2 pts ; question 5 : 2 pts.

\textbf{Points de vigilance :} le calcul doit être présenté avec l'application numérique et l'unité (\%) ; la question 4 exige le mot « prédiction » ou « prévision » ; ne pas écrire que Mendeleïev utilisait déjà le numéro atomique (inconnu à son époque).
\end{corrige}

\begin{corrige}[Exercice 11 — Type BEPC : sécurité domestique et famille des halogènes]
\begin{enumerate}
\item Le gaz dégagé est le \cle{dichlore}, de formule \ce{Cl2} : c'est un gaz verdâtre, irritant et toxique.
\item Le chlore appartient à la famille des \cle{halogènes} (17\textsuperscript{e} colonne). Deux autres éléments de cette famille : le fluor \ce{F} et le brome \ce{Br} (ou l'iode \ce{I}).
\item Le chlore a pour numéro atomique $Z = 17$ ; il se situe dans la \textbf{3\textsuperscript{e} période} du tableau.
\item D'après l'équation \ce{ClO- + Cl- + 2H+ -> Cl2 + H2O}, les ions hypochlorite \ce{ClO-} de l'eau de Javel réagissent avec les ions \ce{H+} du détartrant acide pour former du \textbf{dichlore} \ce{Cl2}, un gaz toxique qui irrite les voies respiratoires et peut provoquer de graves intoxications. Il ne faut donc jamais mélanger ces deux produits.
\item Consigne de sécurité proposée : « Ne mélange jamais l'eau de Javel avec un détartrant ou un autre produit d'entretien : cela dégage un gaz toxique. Range les produits d'entretien fermés, séparés les uns des autres, hors de portée des enfants et loin des aliments. »
\end{enumerate}
\textbf{Barème indicatif (sur 10) :} question 1 : 1,5 pt ; question 2 : 2 pts ; question 3 : 1,5 pt ; question 4 : 3 pts ; question 5 : 2 pts.

\textbf{Points de vigilance :} la formule du gaz doit être \ce{Cl2} (et non \ce{Cl}, qui est le symbole de l'élément) ; la justification de la question 4 doit citer explicitement les réactifs (\ce{ClO-} et \ce{H+}) et le produit toxique (\ce{Cl2}) ; la consigne rédigée doit comporter une interdiction et une mesure de rangement.
\end{corrige}

\newpage

\coursec{Fiche bilan}

\begin{popfigure}
\begin{tikzpicture}[mindmap, grow cyclic, every node/.style={concept}, text width=2.6cm, align=center, root concept/.append style={font=\bfseries, concept color=popblue, text=white, minimum size=3.2cm}, level 1/.append style={level distance=4.3cm, sibling angle=51.4}, level 1 concept/.append style={font=\small, text width=2.9cm}]
\node[root concept]{Classification des éléments}
  child[concept color=poporange, text=popdark]{ node {Élément et symbole} }
  child[concept color=popgreen, text=popdark]{ node {Mendeleïev 1869 : tri par masse et propriétés} }
  child[concept color=poppurple, text=white]{ node {7 lignes = périodes} }
  child[concept color=popblue!70, text=white]{ node {18 colonnes = familles} }
  child[concept color=poppink, text=popdark]{ node {Familles : alcalins, halogènes, gaz rares} }
  child[concept color=popgold, text=popdark]{ node {Métaux / non-métaux} }
  child[concept color=popmuted, text=white]{ node {Éléments usuels et usages} };
\end{tikzpicture}
\legende{Carte mentale de la leçon : les sept idées clés de la classification des éléments chimiques.}
\end{popfigure}

\begin{bilan}
\textbf{Définitions clés}
\begin{itemize}
\item Un \cle{élément chimique} est un type d'atome caractérisé par son numéro atomique $Z$ (nombre de protons).
\item Le \cle{symbole chimique} est une abréviation internationale : une majuscule, parfois suivie d'une minuscule (\ce{H}, \ce{O}, \ce{Na}, \ce{Fe}).
\item Le \cle{tableau périodique} (de Mendeleïev) range les éléments par numéro atomique croissant.
\item Une \cle{période} est une ligne horizontale (7 périodes) ; une \cle{famille} est une colonne verticale (18 colonnes).
\end{itemize}

\textbf{À connaître par cœur}
\begin{center}
\begin{tabularx}{\linewidth}{|l|X|}\hline
\rowcolor{popblueL} \textbf{Notion} & \textbf{Essentiel} \\ \hline
Période & Éléments d'une même ligne : même nombre de couches électroniques. \\ \hline
Famille (colonne) & Éléments de propriétés chimiques voisines (même nombre d'électrons externes). \\ \hline
Alcalins (colonne 1) & \ce{Li}, \ce{Na}, \ce{K} : métaux mous, réagissent violemment avec l'eau. \\ \hline
Halogènes (colonne 17) & \ce{F}, \ce{Cl}, \ce{Br}, \ce{I} : non-métaux très réactifs. \\ \hline
Gaz rares (colonne 18) & \ce{He}, \ce{Ne}, \ce{Ar} : gaz inertes, très stables. \\ \hline
Métaux / non-métaux & Les métaux (majorité) sont à gauche et au centre ; les non-métaux à droite. \\ \hline
\end{tabularx}
\end{center}

\textbf{Méthode express}
\begin{enumerate}
\item Pour situer un élément : repérer sa case, lire le numéro de la ligne (période) et de la colonne (famille).
\item Pour prévoir un comportement chimique : regarder la famille (les éléments d'une même colonne se ressemblent).
\item Pour identifier un symbole : vérifier la majuscule initiale, puis la minuscule éventuelle (\ce{Co} = cobalt, mais \ce{CO} = molécule de monoxyde de carbone !).
\end{enumerate}

\textbf{Éléments usuels au quotidien} : \ce{Fe} (fer, construction), \ce{Al} (aluminium, marmites), \ce{Cu} (cuivre, fils électriques), \ce{Na} et \ce{Cl} (sel de cuisine \ce{NaCl}), \ce{Ca} (calcaire, os), \ce{C} (charbon, essence), \ce{O} et \ce{H} (air, eau), \ce{Au} (or, bijoux).
\end{bilan}

\begin{aretenir}[Checklist avant le BEPC]
\begin{itemize}
\item[$\square$] Je sais définir un élément chimique et donner son symbole correctement (majuscule puis minuscule).
\item[$\square$] Je connais les symboles des éléments usuels : \ce{H}, \ce{O}, \ce{C}, \ce{N}, \ce{Na}, \ce{Cl}, \ce{Fe}, \ce{Cu}, \ce{Al}, \ce{Ca}, \ce{Au}, \ce{Ag}.
\item[$\square$] Je sais qui est Mendeleïev et sur quels critères il a construit son tableau (1869).
\item[$\square$] Je distingue une période (ligne) d'une famille (colonne).
\item[$\square$] Je sais qu'il y a 7 périodes et 18 colonnes.
\item[$\square$] Je cite les trois familles remarquables : alcalins, halogènes, gaz rares, avec un exemple chacune.
\item[$\square$] Je sais que les éléments d'une même famille ont des propriétés chimiques voisines.
\item[$\square$] Je situe les métaux (gauche et centre) et les non-métaux (droite) dans le tableau.
\item[$\square$] Je ne confonds pas \ce{Co} (cobalt) et \ce{CO} (monoxyde de carbone).
\item[$\square$] Je sais associer un élément usuel à une application de la vie courante (cuivre $\to$ câbles ENEO, fer $\to$ ferronnerie, chlore $\to$ eau de Javel).
\end{itemize}
\end{aretenir}

\findecours

\end{document}
